% Template LaTeX Tesis Ciencia de la Computación UCSP
% 2022

\documentclass[a4paper,openany,12pt]{book}
\usepackage[utf8]{inputenc}
\usepackage{graphicx}
\usepackage[spanish,mexico]{babel}
\usepackage{sty/fancyhdr}
\usepackage{ae}
\usepackage[left=2.5cm,right=2.5cm,top=3cm,bottom=2cm]{geometry}
\usepackage[printonlyused]{acronym}
\usepackage{xspace}
\usepackage{sty/hlundef}
\usepackage{sty/tesis}
\usepackage{setspace}
\usepackage{booktabs}
\usepackage{tabularx}

\usepackage{listings}
\usepackage{xcolor}
\usepackage{colortbl}
\definecolor{rowgray}{gray}{0.94}

\lstdefinelanguage{json}{
    basicstyle=\ttfamily\small,
    numbers=left,
    numberstyle=\tiny\color{gray},
    stepnumber=1,
    showstringspaces=false,
    breaklines=true,
    frame=single,
    backgroundcolor=\color{gray!10},
    literate=
     *{0}{{{\color{blue}0}}}{1}
      {1}{{{\color{blue}1}}}{1}
      {2}{{{\color{blue}2}}}{1}
      {3}{{{\color{blue}3}}}{1}
      {4}{{{\color{blue}4}}}{1}
      {5}{{{\color{blue}5}}}{1}
      {6}{{{\color{blue}6}}}{1}
      {7}{{{\color{blue}7}}}{1}
      {8}{{{\color{blue}8}}}{1}
      {9}{{{\color{blue}9}}}{1}
      {:}{{{\color{red}{:}}}}{1}
      {,}{{{\color{red}{,}}}}{1}
      {\{}{{{\color{red}{\{}}}}{1}
      {\}}{{{\color{red}{\}}}}}{1}
      {[}{{{\color{red}{[}}}}{1}
      {]}{{{\color{red}{]}}}}{1},
}


\usepackage{float}
\usepackage[colorlinks=true, citecolor=blue, linkcolor=blue, urlcolor=blue]{hyperref}
\usepackage{comment}
\usepackage[T1]{fontenc}
\usepackage{amssymb} 
\usepackage{array}   
\usepackage{microtype}

\usepackage{longtable}
\usepackage{multirow}
\usepackage{rotating}

\title{** Título de la Tesis **}


%TITULO ANTERIOR
%\title{Sistema de Visual Analytics basado en Grandes Modelos de Lenguaje para la exploración de narrativas complejas con componentes espacio temporales en documentos textuales}

\title{Exploración visual de narrativas espacio temporales en documentos gubernamentales usando Grandes Modelos de Lenguaje}


\author{ Italo Alex Kancha Vilca }


\advisor{Mg. Gina Lucia Muñoz Salas }


\date{Arequipa, Abril 2026}

%\examinerone{Prof. Dr. Hidalgo Buena Gente}{Presidente}%
%\examinertwo{Prof. Dr. Antero A. Gal Oppe}{Secretario}%
%\examinerthree{Prof. Dr. Liu Xiao Ling}{University of ...} % of being the case

\dedicado{Aquí deberás colocar a quien va dedicada tu tesis por ejemplo: A Dios, por todo lo que me ha dado, a todos los profesores por sus enseñanzas y algunos amigos.}

\begin{document}
\pagestyle{fancy}

\maketitle %Compone la carátula y la dedicatoria
\newpage

%\approved{\tres}%  {\tres} or {\cuatro}

%mayores detalles de como usas las abreviaturas (acronimos)
% vea: http://www.ctan.org/tex-archive/macros/latex/contrib/acronym/
% hay un manual en pdf en esa misma direccion

\chapter*{Abreviaturas}

\begin{acronym}
    \acro{LLM}[LLMs]{\textit{Large Language Models} (Grandes Modelos de Lenguaje)}
    \acro{VA}{\textit{Visual Analytics}}
    \acro{NLP}{\textit{Natural Language Processing} (Procesamiento de Lenguaje Natural) }
    \acro{MCV}[MCV]{\textit{Multiple Coordinated Views} (Vistas Múltiples Coordinadas)}
    \acro{OCR}{\textit{Optical Character Recognition} (Reconocimiento Óptico de Caracteres)}
    \acro{PDF}{\textit{Portable Document Format}}
    \acro{JSON}{\textit{JavaScript Object Notation}}
    \acro{API}[APIs]{\textit{Application Programming Interface}}
    \acro{RAG}{\textit{Retrieval-Augmented Generation}}
    \acro{BM25}{\textit{Best Matching 25}}
    \acro{BERT}{\textit{Bidirectional Encoder Representations from Transformers}}
    \acro{SBERT}{\textit{Sentence-BERT}}
    \acro{CRF}{\textit{Conditional Random Field}}
    \acro{LDA}{\textit{Latent Dirichlet Allocation}}
    \acro{MDS}{\textit{Multidimensional Scaling}}
    \acro{PCA}{\textit{Principal Component Analysis}}
    \acro{UMAP}{\textit{Uniform Manifold Approximation and Projection}}
    \acro{TFIDF}[TF-IDF]{\textit{Term Frequency-Inverse Document Frequency}}
    \acro{PSO}{\textit{Particle Swarm Optimization}}
    \acro{SVM}{\textit{Support Vector Machine}}
    \acro{tSNE}[t-SNE]{\textit{t-distributed Stochastic Neighbor Embedding}}

\end{acronym}


\input{Agradecimientos} %Inserta los agradecimientos
\input{Resumen} %Inserta el resumen
\input{Abstract} %Inserta el abstract

\pagenumbering{roman}
\setcounter{page}{1}
\pagestyle{plain}

\tableofcontents %Inserta el índice general
\listoftables %Inserta el índice de cuadros
\listoffigures %Inserta el índice de figuras

%%%%%%%%%%%%%%%%%%%%%%%%%%%%%%%%%%%%%%%%%%%%%%%%%%%%%%%%%%%%%%%%%%%%%
%%%%   En esta parte deberas incluir los archivos de tu tesis   %%%%%
%%%%%%%%%%%%%%%%%%%%%%%%%%%%%%%%%%%%%%%%%%%%%%%%%%%%%%%%%%%%%%%%%%%%%


\pagestyle{plain}
\pagenumbering{arabic}
\setcounter{page}{1}
% TO DO
% Italizar
% Revisar nuevamente las fuentes 
% eso de paradigma de VA no m termina de convencer
% Aporte computacional en la introducción — reforzar que las preguntas de investigación solo pueden responderse construyendo y evaluando el sistema, para que los objetivos se lean como investigación y no como ingeniería.
%Poner bonito la organizacion d la tesis
\chapter{Introducción}
\label{cap:introduccion}

Cada año, instituciones públicas como tribunales, organismos de 
control y agencias regulatorias producen miles de documentos que 
registran el avance de procesos que pueden extenderse durante meses 
o años. Identificar qué ocurrió en un caso, cuándo y dónde, sigue 
requiriendo que un analista lea cientos de documentos de forma 
manual, porque la información está dispersa en archivos generados 
por distintos actores en distintos momentos. Reconstruir 
automáticamente la secuencia de eventos que conforma cada caso 
sigue siendo un problema abierto. Los sistemas tradicionales de 
extracción de información operan a nivel de documento individual y 
los enfoques basados en reglas requieren plantillas específicas por 
tipo de documento, lo que dificulta integrar información cuando una 
misma entidad aparece referenciada de formas distintas en archivos 
generados por distintos actores~\cite{cattan2021crossdoc}.

Los Grandes Modelos de Lenguaje (\acs{LLM}, por sus siglas en inglés)
 han demostrado capacidades 
prometedoras para extraer entidades, relaciones y secuencias 
narrativas de documentos técnicos con una flexibilidad que los 
enfoques basados en reglas no alcanzan~\cite{brown2020, wei2022}. 
Sin embargo, en tareas de extracción estructurada estos modelos 
presentan tasas de alucinación significativas~\cite{ji2023}, lo 
cual representa un riesgo inaceptable cuando una fecha incorrecta 
o un monto fabricado puede comprometer una decisión institucional. 
Si bien parte de estas inconsistencias puede detectarse 
automáticamente contrastando cada dato contra el texto fuente, la 
corrección automática no resuelve el problema de fondo. El analista 
necesita un marco que le permita calibrar su confianza en lo 
extraído mientras explora cómo evolucionó el proceso. El 
\acl{VA} (\acs{VA}) ~\cite{keim2008visual, sacha2014} 
ofrece ese marco. A diferencia de las interfaces de consulta 
basadas en chat, que devuelven respuestas puntuales sobre un 
documento a la vez, %Creo que aqui falta una cita :)
\ac{VA} permite mantener simultáneamente las 
dimensiones temporal, geoespacial y semántica de un proceso 
distribuido en múltiples documentos, integrando el procesamiento 
automático con el razonamiento humano. Esta integración aún no ha 
sido evaluada en documentos institucionales fragmentados en español.

En este trabajo abordamos esa brecha mediante el diseño y 
evaluación de un sistema de \ac{VA} que articula tres componentes: un 
pipeline de extracción basado en \acs{LLM} con corrección automática de 
inconsistencias, un mecanismo de trazabilidad que vincula cada 
dato con su fragmento de origen, y un conjunto de vistas 
interactivas coordinadas para las dimensiones temporal, geoespacial 
y narrativa de los procesos. Concretamente, esta investigación 
responde dos preguntas. \textbf{(1)} ¿Cómo pueden los \acs{LLM} automatizar la extracción de información estructurada a partir de documentos institucionales fragmentados en español? \textbf{(2)} 
¿Qué formas de visualización e interacción permiten a un 
analista explorar y verificar narrativas complejas extraídas 
de documentos institucionales fragmentados? Para 
responderlas, evaluamos el sistema con expedientes reales de 
auditoría de la Contraloría General de la República del Perú, un contexto real con alto volumen documental en español.




\section{Motivación y Contexto}

Los organismos de control gubernamental generan grandes 
volúmenes de documentos que registran procesos extendidos en 
el tiempo. En el Perú, la Contraloría General de la República 
procesó más de 50,000 informes solo en 
2024~\cite{andina2024}. Estos documentos presentan 
características que los convierten en un caso de estudio 
relevante para la investigación en análisis visual de texto: 
combinan narrativa densa con terminología especializada, 
involucran múltiples actores y ubicaciones geográficas, y 
registran secuencias de eventos que se extienden durante 
meses a través de documentos generados de forma 
independiente.

Si bien el \ac{NLP} permite clasificar y extraer elementos 
clave de documentos legales~\cite{zhong-etal-2020-nlp, 
lippi2019legal}, estos enfoques operan predominantemente a 
nivel de documento individual y requieren adaptación 
específica por dominio. Los \acs{LLM}, en cambio, permiten 
adaptar la extracción a nuevos tipos de documentos sin 
reentrenamiento, lo que los convierte en una alternativa 
de propósito general para múltiples dominios.

El \ac{VA} aborda este tipo de problemas combinando procesamiento automático con representación visual interactiva, permitiendo que el analista explore grandes volúmenes de información mientras mantiene control sobre el proceso de interpretación~\cite{keim2008visual, sacha2014}. Esta combinación es especialmente relevante cuando los datos extraídos automáticamente requieren validación experta antes de sustentar decisiones.

Sin embargo, reconstruir la secuencia completa de un proceso 
a partir de múltiples documentos fragmentados requiere 
resolver desafíos adicionales tales como identificar que 
distintas menciones en distintos archivos se refieren a la 
misma entidad~\cite{cattan2021crossdoc}, ordenar 
temporalmente eventos descritos con formatos heterogéneos, 
y consolidar información dispersa en una representación 
coherente. Permitir además que el analista verifique esa 
reconstrucción sigue siendo un problema 
abierto~\cite{pereira2023inacia}. Abordar este problema de 
forma integrada, combinando extracción automática con 
representación visual interactiva, es el punto de partida 
de esta investigación.

\section{Planteamiento del Problema}

El análisis de procesos documentados en texto no estructurado 
plantea tres problemas computacionales sin resolver.

El primero es la extracción y resolución de entidades a través 
de múltiples documentos. Los sistemas existentes operan 
predominantemente a nivel de documento individual y no resuelven 
la correferencia entre documentos: la misma entidad puede aparecer 
referenciada de formas distintas en distintos archivos, y sin un 
mecanismo de resolución de identidad los datos extraídos no pueden 
integrarse de manera coherente. Los \acs{LLM} ofrecen capacidades 
prometedoras para este problema, pero su comportamiento en 
documentos técnicos y jurídicos en español no ha sido 
sistemáticamente evaluado en pipelines de extracción con revisión 
asistida~\cite{zhong-etal-2020-nlp}.

El segundo es la verificabilidad de la extracción automática. 
Cuando un sistema extrae un dato, el analista no dispone de 
mecanismos inmediatos para verificar su corrección sin consultar 
manualmente el documento fuente. En dominios donde un dato incorrecto puede sustentar una decisión institucional, esta ausencia de verificación hace que los sistemas automáticos sean difíciles de adoptar en la práctica~\cite{ji2023}.

El tercero es la representación visual integrada de procesos 
multidimensionales. Los procesos documentados en texto tienen 
simultáneamente una dimensión temporal, espacial y narrativa. 
Caracterizar qué representaciones visuales e interacciones permiten explorar de forma integrada las dimensiones temporal, espacial y narrativa de un proceso distribuido en múltiples documentos sigue siendo un problema abierto~\cite{zhao2023contextwing, rajabiyazdi2024textvista, ye2024storyexplorer, yeh2025storyribbons}.
\section{Objetivos}

\subsection{Objetivo General}

Desarrollar un sistema de \acl{VA} con extracción automática de 
información mediante \acs{LLM} y revisión asistida, para la 
exploración interactiva de procesos documentados en texto 
no estructurado a través de sus dimensiones temporal, 
geoespacial y narrativa.

\subsection{Objetivos Específicos}

\begin{enumerate}
    \item Diseñar un \textit{pipeline} de extracción basado en \acs{LLM} 
     que identifique entidades, hitos temporales, actores, 
     relaciones y ubicaciones en documentos textuales, 
     mitigando posibles alucinaciones del modelo.

     % ambas vistas son formas complementarias de explorar los datos extraídos, por eso están en el mismo objetivo
     \item Desarrollar un prototipo de vistas interactivas 
     coordinadas que permitan explorar las dimensiones 
     temporal, geoespacial y narrativa de los procesos 
     documentados, incluyendo una vista de similitud 
     semántica que permita comparar entidades entre sí.
     
     %\item Desarrollar un mecanismo de trazabilidad que preserve la referencia al documento fuente durante todo el pipeline, permitiendo al analista acceder al fragmento original de cada dato desde la interfaz.
     
     \item Evaluar la precisión del pipeline de extracción 
     mediante un caso de estudio con expedientes de la 
     Contraloría General de la República del Perú, 
     comparando las extracciones generadas contra el 
     contenido verificado de los documentos fuente.
     
     \item Evaluar la utilidad y usabilidad del sistema 
     mediante una prueba con usuarios, midiendo eficiencia 
     en tareas de exploración documental y utilidad 
     percibida mediante cuestionarios estandarizados.
\end{enumerate}


\section{Organización de la tesis}

El presente documento se organiza de la siguiente manera:

\begin{itemize}
    \item \hyperref[cap:introduccion]{\textbf{Capítulo~\ref*{cap:introduccion}: Introducción}} presenta el problema de investigación, la motivación, las preguntas de investigación y los objetivos del trabajo.
    \item \hyperref[cap:marco_teorico]{\textbf{Capítulo~\ref*{cap:marco_teorico}: Marco Teórico}} desarrolla el marco teórico, 
     cubriendo los fundamentos de \ac{VA}, visualización 
     espacio-temporal, extracción de información con \acs{LLM} 
      y el dominio de Control Concurrente de la Contraloría 
      General de la República del Perú.
     \item \hyperref[cap:estado_arte]{\textbf{Capítulo~\ref*{cap:estado_arte}: Estado del Arte}} presenta el estado del arte, 
      revisando los trabajos relacionados en extracción de 
      información con \acs{LLM}, sistemas de \ac{VA} para texto y 
      visualización de narrativas complejas.
      \item \hyperref[cap:propuesta]{\textbf{Capítulo~\ref*{cap:propuesta}: Propuesta}} describe el diseño e implementación 
      del sistema, incluyendo el pipeline de extracción, la 
      corrección automática de inconsistencias y las vistas 
      interactivas coordinadas.
      \item El Capítulo 5 presenta la evaluación del sistema 
      en tres partes: la precisión del pipeline y el impacto 
      de la corrección automática sobre la tasa de 
      alucinaciones, la calidad de la proyección semántica 
      para identificación de patrones, y la prueba de 
      usabilidad con usuarios.
      \item El Capítulo 6 discute las conclusiones, las 
      limitaciones del trabajo y las líneas de investigación 
      futura.
\end{itemize}

\section{Cronograma}

\begin{figure}[H]
    \centering
    \includegraphics[width=0.9\textwidth]{Tesis_CS_202203/Cro_Actv.pdf}
    \caption{Cronograma de actividades.}
    \label{fig:cronograma}
\end{figure}
 %Inserta el capítulo 1
%TODO
%Italizar
%revisar las referencias
%revisar todo cuando acabe la propuesta
%cambiar la imagen delas vistsas coordinadas
%actualizar consideraciones iniciales
%mas vis d tiempo
%interaccion
%explayarte en reducciones dimensionales
%Sacar vector caracteristicas
%LLMS escueto

\chapter{Marco Teórico}
\label{cap:marco_teorico}


\section{Consideraciones Iniciales}
Este capítulo presenta los conceptos teóricos que 
fundamentan la propuesta. Se comienza con \ac{VA}, 
su ciclo iterativo y el patrón de vistas coordinadas 
que organiza la interfaz. Luego, se revisa la 
visualización de procesos con dimensiones temporales 
y espaciales, y cómo la trazabilidad permite verificar 
los datos extraídos. Después se describen los 
\acs{LLM}, su uso para extraer información estructurada 
de texto, los problemas que esto implica y estrategias de mitigación. El capítulo 
cierra presentando los documentos institucionales 
fragmentados y el corpus 
utilizado para evaluar el sistema.

\section{\textit{Visual Analytics}}

\ac{VA} se define como la ciencia del razonamiento 
analítico facilitado por interfaces visuales 
interactivas~\cite{thomas2005illuminating}. Su núcleo 
es la integración del procesamiento automático con las 
capacidades de razonamiento humano, de modo que el 
analista participa activamente en el proceso de 
análisis y no solo observa sus resultados~\cite{keim2008visual, 
cui2019visual}.

A partir del modelo de referencia de Keim ~\cite{keim2008visual}, el proceso opera como un 
ciclo iterativo entre cuatro etapas interdependientes, 
tal como se ilustra en la Figura~\ref{fig:keim_va}. 
Primero, los datos crudos se transforman en estructuras 
manejables mediante tareas de limpieza y normalización. 
Luego toman dos vías simultáneas: la construcción de 
modelos automáticos que extraen patrones, y el mapeo de 
los resultados en representaciones visuales interactivas. 
A partir de esas representaciones, el analista evalúa 
los resultados, identifica patrones o anomalías y formula 
nuevas hipótesis. El 
conocimiento adquirido se reintegra al sistema mediante 
ajustes en los parámetros de los modelos o en los filtros 
de las vistas, reiniciando el ciclo.

\begin{figure}[H]
\centering
\includegraphics[width=\textwidth]{figs/CicloVA.pdf}
\caption{Ciclo de \ac{VA}, mostrando la interacción continua 
entre los datos, los modelos automáticos, las representaciones 
visuales y el analista. Adaptado de~\cite{keim2008visual}.}
\label{fig:keim_va}
\end{figure}

Sacha et al.~\cite{sacha2014} refinan este modelo 
distinguiendo explícitamente dos componentes que operan 
de forma complementaria. 
El componente computacional comprende los datos, los 
modelos analíticos y las visualizaciones generadas a 
partir de ellos. El componente humano comprende el 
proceso cognitivo del analista, que se estructura en 
tres bucles de razonamiento anidados.

El \textbf{bucle de exploración} (\textit{exploration loop}) 
describe cómo el analista interactúa con el sistema 
mediante acciones que producen hallazgos. Una acción 
puede ser filtrar un subconjunto de datos, cambiar el 
mapeo visual o ajustar los parámetros de un modelo. El 
sistema responde a esa acción y el analista observa el 
resultado, generando un hallazgo (\textit{finding}). Un 
hallazgo es cualquier observación relevante que el 
analista percibe en la visualización o en el modelo, 
como un patrón inesperado, un valor atípico o la 
ausencia de un dato esperado. Este bucle se repite de 
forma rápida y constituye la base de toda exploración 
interactiva.

El \textbf{bucle de verificación} (\textit{verification loop}) 
opera a un nivel superior. Cuando el analista interpreta 
un hallazgo en el contexto del dominio del problema, 
produce un \textit{insight}, una unidad de comprensión 
que conecta la evidencia visual con el conocimiento 
previo. Los \textit{insights} pueden confirmar o refutar 
una hipótesis existente, o generar hipótesis nuevas que 
requieran exploración adicional. Este bucle vincula lo 
que el analista observa en la interfaz con lo que sabe 
del dominio.

El \textbf{bucle de generación de conocimiento} 
(\textit{knowledge generation loop}) cierra el proceso. 
Cuando el analista acumula suficiente evidencia para 
confiar en un \textit{insight}, este se consolida como 
conocimiento. Ese conocimiento nuevo puede a su vez 
generar hipótesis adicionales, reiniciando el ciclo 
completo. Sacha et al. enfatizan que el conocimiento 
no se produce en el componente computacional ni en el 
humano por separado, sino en la interacción entre 
ambos~\cite{sacha2014}. El analista aporta juicio 
contextual y razonamiento que los algoritmos no poseen, 
mientras que los algoritmos procesan volúmenes de 
información que el analista no podría revisar 
manualmente. Esta complementariedad es la que 
permite que \ac{VA} funcione como un marco de generación 
de conocimiento.



La distinción entre estos tres bucles resulta relevante 
para el diseño de sistemas de \ac{VA} porque permite 
evaluar qué niveles de razonamiento soporta cada 
componente. Un sistema puede facilitar el 
bucle de exploración ofreciendo visualizaciones 
interactivas, pero si carece de mecanismos para que el 
analista verifique la evidencia, el bucle de verificación 
queda sin soporte. 





\subsection{\textit{Visual Analytics} aplicado a documentos 
textuales}
\label{sec:va_texto}

El procesamiento de información textual representa uno 
de los mayores desafíos dentro de \ac{VA} debido a la 
naturaleza no estructurada del lenguaje. 
Thomas y Cook  ~\cite{thomas2005illuminating} señalan que 
los documentos en texto libre requieren transformaciones 
semánticas profundas para extraer unidades de significado 
antes de poder ser representados visualmente o analizados 
algorítmicamente. A diferencia de datos numéricos o 
tabulares, donde las dimensiones del análisis están 
definidas de antemano, en un documento textual las 
entidades, las relaciones y la estructura temporal deben 
ser identificadas y extraídas como paso previo a 
cualquier representación visual. 
Kucher y Kerren~\cite{kucher2015text} proponen una 
taxonomía de técnicas de visualización de texto que 
organiza este campo según la tarea analítica, el tipo 
de dato y las propiedades de la representación visual. 
Esta taxonomía permite situar cualquier sistema de 
\ac{VA} para texto dentro de un espacio de diseño 
definido por las decisiones que conectan el dato textual 
con su forma visual.

Cuando los documentos provienen de dominios donde las 
decisiones tienen consecuencias institucionales, la 
extracción automática enfrenta una exigencia adicional. 
Los algoritmos que operan como cajas negras dificultan 
que el experto humano pueda auditar y confiar en los 
resultados. Hutchinson et al.~\cite{hutchinson2024llm} 
advierten que para incorporar modelos de lenguaje dentro 
de un entorno de \ac{VA}, sus respuestas deben estar 
ancladas a los datos originales. Esta necesidad de 
anclaje conecta directamente con el bucle de verificación 
de Sacha et al.~\cite{sacha2014}, donde el analista 
evalúa críticamente los hallazgos producidos por la 
máquina para determinar si constituyen evidencia real o 
un artefacto del procesamiento automático.

\subsubsection{Trazabilidad y lectura \textit{in situ}}
\label{sec:trazabilidad}

La trazabilidad es el mecanismo que vincula cada 
elemento visual de la interfaz con el fragmento 
específico del documento del que fue extraído. En 
términos del modelo de Sacha et al.~\cite{sacha2014}, 
este mecanismo soporta el bucle de verificación al 
permitir que el analista pase de un hallazgo visual a 
la evidencia textual que lo sustenta. Sin este vínculo, 
el analista solo puede aceptar o rechazar los resultados 
del modelo sin poder evaluar su corrección.

La lectura \textit{in situ} es la forma concreta en que 
se materializa la trazabilidad en la interfaz. Consiste 
en permitir que el analista, al seleccionar un dato en 
cualquiera de las vistas, acceda al fragmento exacto 
del documento original del que fue extraído, incluyendo 
su contexto circundante y su ubicación en la página. 
Este mecanismo permite un flujo de trabajo en el que 
el analista transita entre la vista panorámica de los 
datos y la lectura detallada del texto fuente sin 
abandonar la interfaz~\cite{thomas2005illuminating}.







\section{Visualización de Datos}
\label{sec:tecnicas_vis}

La extracción automática de información resulta 
insuficiente si el usuario no puede interpretar los 
resultados. La visualización de datos actúa como el 
canal que transforma información abstracta en 
representaciones gráficas, aprovechando la capacidad 
perceptiva humana para revelar patrones que de otro 
modo permanecerían ocultos en los 
datos~\cite{ortigossa2025time}. Esta sección presenta 
los principios de interacción, los patrones de 
coordinación entre vistas y las técnicas de 
representación que fundamentan el diseño de la 
interfaz.


\subsection{El Mantra de Shneiderman}
\label{sec:shneiderman}

El diseño de interfaces exploratorias se fundamenta en 
el principio propuesto por 
Shneiderman~\cite{shneiderman1996eyes}, conocido como 
el Mantra de Búsqueda de Información Visual. Este 
principio organiza la exploración en tres etapas 
jerárquicas. En la primera, \textit{overview}, el 
sistema presenta una perspectiva global del conjunto 
de datos que permite al analista orientarse sin 
saturación visual. En la segunda, \textit{zoom and 
filter}, el analista aísla subconjuntos de interés 
mediante controles interactivos, eliminando 
dinámicamente la información irrelevante. En la 
tercera, \textit{details-on-demand}, el analista 
selecciona un elemento específico y accede a su 
información completa. Este flujo descendente, de lo 
general a lo particular, permite que el analista 
construya progresivamente su comprensión del conjunto 
de datos sin necesidad de examinar cada elemento 
individualmente.


\subsection{Vistas múltiples coordinadas}
\label{sec:mcv}

Cuando los datos poseen múltiples dimensiones, una sola 
representación visual rara vez es suficiente para 
explorarlos por completo. Las vistas múltiples coordinadas (\acs{MCV}, pos sus siglas en ingles) son un patrón de 
diseño en el que dos o más representaciones visuales del 
mismo conjunto de datos se presentan de forma simultánea, 
y las acciones del usuario en una vista se propagan 
automáticamente a las 
demás~\cite{roberts2007state, baldonado2000guidelines}. 
De este modo, cada vista puede especializarse en una 
dimensión del análisis mientras la coordinación entre 
ellas preserva el contexto compartido 
(Figura~\ref{fig:brushing_linking}).

El mecanismo fundamental de coordinación es la selección y vinculación interactiva 
(\textit{brushing and linking})~\cite{becker1987brushing}. 
Cuando el analista selecciona un elemento en una vista, 
esa selección se transmite a todas las vistas vinculadas, 
donde el mismo elemento queda resaltado en su 
representación correspondiente. Un ejemplo de esto puede observarse en la 
(Figura~\ref{fig:brushing_linking}). 
Becker y Cleveland~\cite{becker1987brushing} definen 
cuatro operaciones de \textit{brushing} (resaltado, 
resaltado con sombra, eliminación y etiquetado) cuyo 
efecto se propaga simultáneamente a todas las vistas. 
Este mecanismo permite que el analista observe cómo se 
manifiesta una misma entidad en distintas dimensiones 
sin perder la referencia de qué elemento está analizando.

% FIGURA: Adaptar de Becker & Cleveland 1987, Figure 1
% Muestra scatterplots con brushing: seleccionar puntos 
% en un scatterplot los resalta en los demás.
% También puedes adaptar de Roberts 2007, Figure 1 
% (Improvise) que muestra 6 vistas coordinadas.
\begin{figure}[H]
\centering
\includegraphics[width=0.85\textwidth]%
{figs/improvise.png}
\caption{Sistema Improvise mostrando seis vistas 
coordinadas del mismo conjunto de datos. La selección 
de elementos en una vista se propaga a las demás 
mediante \textit{brushing and linking}. Imagen 
extraída de~\cite{roberts2007state}.}
\label{fig:brushing_linking}
\end{figure}

Además del \textit{brushing}, 
Roberts~\cite{roberts2007state} identifica otros 
mecanismos de coordinación entre vistas. La navegación 
sincronizada permite que un cambio de zoom o 
desplazamiento en una vista se replique en las demás. 
El filtrado compartido permite que los criterios de 
exclusión aplicados en una vista reduzcan el conjunto 
visible en todas las vistas vinculadas. La combinación 
de estos mecanismos permite al analista explorar 
distintas facetas de los datos de forma fluida, 
manteniendo la coherencia entre las representaciones.

Desde una perspectiva computacional, la coordinación 
entre vistas implica mantener un estado compartido de 
selección y filtrado que se actualiza y propaga cada 
vez que el usuario interactúa con cualquiera de las 
vistas. 
%North y Shneiderman~\cite{north2000snap}  formalizan este mecanismo utilizando el modelo relacional de datos, donde las relaciones entre registros de distintas tablas determinan qué elementos se vinculan entre vistas. Esta formalización permite definir la coordinación de forma independiente de las representaciones visuales específicas que se utilicen.

Baldonado et al.~\cite{baldonado2000guidelines} 
proponen directrices para el diseño de estos sistemas. 
Entre ellas destacan el principio de diversidad, 
que establece que las vistas deben usar representaciones 
distintas para aportar perspectivas complementarias 
sobre los mismos datos, y el principio de 
parsimonia, que recomienda usar el menor número 
de vistas necesario para cubrir las dimensiones del 
análisis sin imponer una carga cognitiva excesiva al 
usuario. Estos principios implican que cada vista 
adicional debe justificarse por la perspectiva nueva 
que aporta, y que la coordinación entre ellas debe 
reducir, y no aumentar, el esfuerzo cognitivo del 
analista.

En el marco de \ac{VA}, las \acs{MCV} cumplen un rol 
especialmente importante porque permiten sostener los 
bucles de razonamiento descritos por 
Sacha et al.~\cite{sacha2014} a través de múltiples 
dimensiones de forma simultánea. El bucle de 
exploración se beneficia porque el analista puede 
generar hallazgos comparando lo que observa en una 
vista con lo que aparece en otra. El bucle de 
verificación se beneficia porque un dato sospechoso 
en una vista puede contrastarse inmediatamente con 
su representación en una vista diferente.

\subsection{Representación de datos temporales}
\label{sec:temporal}

El tiempo posee una estructura semántica que lo 
distingue de otras variables cuantitativas. 
Aigner et al.~\cite{aigner2011visualization} proponen 
una taxonomía que caracteriza los datos temporales 
según tres propiedades. La escala define si el tiempo 
se representa de forma ordinal (solo el orden importa) 
o cuantitativa (las distancias entre instantes son 
medibles). El alcance distingue entre datos puntuales, 
que registran un instante, y datos basados en 
intervalos, que poseen inicio y fin. El arreglo 
describe si la progresión temporal es lineal o cíclica. 
Estas propiedades determinan qué técnicas de 
representación son adecuadas para cada tipo de dato 
temporal. Ortigossa et al.~\cite{ortigossa2025time} 
ofrecen una revisión amplia de técnicas de 
visualización para series temporales, organizadas 
según estas propiedades y según el tipo de tarea 
analítica que soportan.

\subsubsection{Gráficos de líneas}
\label{sec:linechart}

El gráfico de líneas es la forma más antigua de 
representación de series temporales. 
Playfair~\cite{playfair1786atlas} introdujo esta 
técnica en 1786 para representar la evolución de las 
exportaciones e importaciones de Inglaterra a lo largo 
del siglo XVIII, mapeando el tiempo al eje horizontal 
y los valores al eje vertical. Dos siglos después, 
esta representación permanece esencialmente sin 
cambios~\cite{tufte2001visual}. Cuando se superponen 
múltiples series en el mismo espacio, el analista puede 
comparar tendencias, identificar puntos de cruce y 
detectar divergencias en la evolución de distintas 
variables. Cleveland~\cite{cleveland1993visualizing} 
establece principios perceptivos para la lectura 
efectiva de estos gráficos, como la elección de la 
relación de aspecto (\textit{aspect ratio}) para 
facilitar la percepción de las pendientes. La principal 
limitación del gráfico de líneas aparece cuando el 
número de series crece, ya que la superposición 
dificulta la lectura individual de cada una. La 
Figura~\ref{fig:line_chart} ilustra la superposición 
de tres series temporales sobre un eje común.

\begin{figure}[H]
\centering
\includegraphics[width=0.8\textwidth]{figs/line_chart.png}
\caption{Gráfico de líneas mostrando la evolución 
temporal de tres variables a lo largo de doce meses. 
La superposición de series permite comparar tendencias 
y detectar divergencias.}
\label{fig:line_chart}
\end{figure}

\subsubsection{Diagramas de Gantt}
\label{sec:gantt}

Para datos basados en intervalos con arreglo lineal, 
los diagramas de Gantt constituyen la representación 
visual más adecuada~\cite{aigner2011visualization, 
ortigossa2025time}. Desarrollado originalmente por 
Henry Gantt a inicios del siglo 
XX~\cite{wilson2003gantt}, este diagrama asocia cada 
entidad a una fila horizontal y representa los 
intervalos mediante barras cuya longitud es proporcional 
a su duración. Esta técnica permite comparar 
visualmente solapamientos entre fases, identificar 
retrasos y evaluar la progresión relativa de múltiples 
procesos sobre un eje temporal 
común~\cite{aigner2011visualization}. La 
Figura~\ref{fig:gantt_time} muestra cinco entidades 
con intervalos de distinta duración y solapamiento.

\begin{figure}[H]
\centering
\includegraphics[width=0.85\textwidth]{figs/gantt_chart.png}
\caption{Diagrama de Gantt mostrando los intervalos 
temporales de cinco entidades. La longitud de cada 
barra es proporcional a la duración del intervalo, 
y la disposición en filas permite identificar 
solapamientos y retrasos entre procesos paralelos. }
\label{fig:gantt_time}
\end{figure}

\subsubsection{\textit{Streamgraph}}
\label{sec:streamgraph}

El \textit{streamgraph} es una variante del gráfico de 
áreas apiladas en la que múltiples series temporales se 
representan como flujos cuyo grosor varía en función del 
valor de cada variable en cada punto del tiempo. 
Havre et al.~\cite{havre2002themeriver} introdujeron 
esta técnica con el nombre de ThemeRiver para 
visualizar la evolución temática de grandes colecciones 
de documentos. Cada tema se representa como una 
corriente cuyo ancho refleja su frecuencia, y el 
conjunto de corrientes forma un flujo continuo a lo 
largo del eje temporal. Byron y 
Wattenberg~\cite{byron2008stacked} formalizaron la 
geometría del \textit{streamgraph}, proponiendo 
algoritmos para optimizar la disposición de las capas 
y minimizar la distorsión visual que surge al apilar 
múltiples series.

A diferencia del gráfico de líneas, que enfatiza los 
valores individuales de cada serie, el 
\textit{streamgraph} enfatiza las proporciones 
relativas y la composición del total a lo largo del 
tiempo~\cite{byron2008stacked}. Esto lo hace 
especialmente útil cuando el analista necesita percibir 
cómo cambia la distribución de un conjunto de variables 
entre períodos. La Figura~\ref{fig:streamgraph} muestra 
cómo cuatro variables compiten en proporción a lo 
largo del tiempo.

\begin{figure}[H]
\centering
\includegraphics[width=0.85\textwidth]{figs/streamgraph.png}
\caption{\textit{Streamgraph} mostrando la evolución 
de cuatro variables a lo largo del tiempo. El grosor 
de cada flujo refleja el valor de la variable en cada 
punto, y la disposición simétrica facilita la 
percepción de las proporciones relativas entre 
variables.}
\label{fig:streamgraph}
\end{figure}

\subsubsection{Diagramas de Sankey}
\label{sec:sankey}

Mientras que el \textit{streamgraph} muestra cómo 
cambia la composición de un conjunto a lo largo del 
tiempo, el diagrama de Sankey muestra cómo fluyen 
cantidades entre estados o categorías. Originalmente 
concebido por el capitán irlandés Matthew 
Sankey en 1898 para representar las pérdidas 
energéticas de una máquina de 
vapor~\cite{schmidt2008sankey}, este tipo de diagrama 
utiliza nodos para representar estados y conexiones 
cuyo ancho es proporcional a la cantidad que transita 
entre ellos. Riehmann et al.~\cite{riehmann2005sankey} 
proponen un sistema interactivo que soporta el flujo 
de exploración de Shneiderman, ofreciendo una vista 
general del grafo de flujos con la posibilidad de 
hacer zoom y acceder a detalles bajo demanda.

Cuando se incorpora la dimensión temporal, los nodos 
se organizan en columnas que representan momentos 
sucesivos, y los flujos entre columnas muestran las 
transiciones entre estados de un momento al siguiente. 
Esta disposición permite identificar qué transiciones 
son frecuentes, dónde se concentran las pérdidas y 
qué trayectorias siguen las entidades a lo largo del 
proceso. La Figura~\ref{fig:sankey} ilustra este 
esquema con tres estados distribuidos en tres momentos 
sucesivos.

\begin{figure}[H]
\centering
\includegraphics[width=0.85\textwidth]{Tesis_CS_202203/figs/sankey_horizontal.png}
\caption{Diagrama de Sankey temporal con tres estados 
(A, B, C) en tres momentos sucesivos ($t_1$, $t_2$, 
$t_3$). El ancho de cada flujo es proporcional a la 
cantidad de entidades que transitan entre estados, 
permitiendo identificar las trayectorias más 
frecuentes.}
\label{fig:sankey}
\end{figure}

\subsection{Representación geoespacial}
\label{sec:geoespacial}

Cuando los datos poseen una dimensión geográfica, las 
técnicas de geovisualización permiten revelar patrones 
espaciales que permanecen invisibles en representaciones 
tabulares. Andrienko y Andrienko~\cite{andrienko2006exploratory} 
establecen que el análisis exploratorio de datos 
espaciales requiere técnicas capaces de representar 
simultáneamente la distribución geográfica de los datos 
y sus atributos temáticos, de modo que el analista pueda 
percibir tanto dónde ocurren los fenómenos como cuál es 
su magnitud.

Las dos técnicas fundamentales de la cartografía temática 
son los mapas coropléticos y los mapas de símbolos 
proporcionales~\cite{slocum2014thematic}. Los mapas 
coropléticos codifican valores agregados mediante la 
variación de color o intensidad en regiones predefinidas, 
lo que los hace adecuados para representar densidades, 
tasas o proporciones cuya distribución sigue límites 
administrativos. Los mapas de símbolos proporcionales 
representan magnitudes mediante marcadores cuyo tamaño 
varía según el valor asociado a cada ubicación, lo que 
los hace preferibles cuando los datos corresponden a 
conteos o totales asociados a puntos específicos. Ambas 
técnicas pueden combinarse en un mismo mapa: las regiones 
aportan el contexto geográfico mediante color, mientras 
que los símbolos superpuestos comunican magnitudes 
puntuales~\cite{slocum2014thematic}. A este tipo de 
representación, que integra polígonos coloreados, 
símbolos y etiquetas en una sola vista, se le denomina 
mapa anotado.

Cuando las entidades se distribuyen en múltiples niveles 
administrativos, la agregación jerárquica permite 
transitar entre escalas. En el nivel más amplio el mapa 
muestra valores agregados por región. Al hacer zoom, los 
agregados se descomponen en las unidades que los conforman, 
revelando la distribución interna. Este mecanismo, conocido 
como \textit{zoom semántico}, ajusta el nivel de detalle 
de la representación según la escala de 
visualización~\cite{andrienko2006exploratory}. La 
Figura~\ref{fig:geo_zoom} ilustra este esquema en tres 
niveles: nacional, regional y local.

% FIGURA: Tres paneles horizontales mostrando zoom semántico
% Panel 1 (nacional): mapa con regiones coloreadas
% Panel 2 (regional): zoom a una región con símbolos proporcionales
% Panel 3 (local): zoom a una ubicación con anotaciones de detalle
\begin{figure}[H]
\centering
\includegraphics[width=\textwidth]{figs/geo_zoom.png}
\caption{Zoom semántico en tres niveles sobre un mapa 
anotado. En el nivel nacional las regiones se colorean 
según un valor agregado. Al hacer zoom a una región 
aparecen símbolos proporcionales por entidad. En el 
nivel local se muestran las anotaciones de detalle de 
una entidad específica.}
\label{fig:geo_zoom}
\end{figure}

La combinación del \textit{zoom semántico} con los 
principios de Shneiderman permite implementar el flujo de \textit{overview}, 
\textit{zoom and filter} y \textit{details-on-demand} 
en la dimensión espacial. El nivel nacional corresponde 
al \textit{overview}; la selección de una región activa 
el \textit{zoom and filter}; y la selección de una 
entidad específica despliega sus \textit{details-on-demand}. 
Esta correspondencia entre los principios de interacción 
y los niveles geográficos convierte al mapa anotado en 
un punto de entrada natural para la exploración de 
datos distribuidos geográficamente.

\subsection{Reducción dimensional y proyección semántica}
\label{sec:reduccion}

Cuando cada entidad de un conjunto de datos se 
describe mediante un vector de alta dimensión, la 
comparación directa entre entidades resulta 
cognitivamente inviable. Las técnicas de reducción 
dimensional proyectan estos vectores en un espacio de 
dos o tres dimensiones preservando, en la medida de lo 
posible, las relaciones de proximidad del espacio 
original. En la proyección resultante, las entidades 
que son similares en el espacio de alta dimensión 
aparecen cercanas en el plano, permitiendo al analista 
identificar agrupaciones, valores atípicos y patrones 
de similitud de forma visual.

t-SNE (\textit{t-distributed Stochastic Neighbor 
Embedding})~\cite{vandermaaten2008tsne} fue una de las 
primeras técnicas en lograr proyecciones de alta calidad 
para visualización, modelando las similitudes entre 
puntos como distribuciones de probabilidad y 
minimizando la divergencia entre las distribuciones 
del espacio original y el proyectado. Sin embargo, 
t-SNE presenta limitaciones en la preservación de la 
estructura global de los datos y en su escalabilidad 
a conjuntos de datos grandes.

\acs{UMAP} (\textit{Uniform Manifold Approximation and 
Projection})~\cite{mcinnes2018umap} aborda estas 
limitaciones. Fundamentado en geometría riemanniana 
y topología algebraica, \acs{UMAP} construye una 
representación topológica difusa del espacio de alta 
dimensión y optimiza su proyección en dimensiones 
bajas minimizando la entropía cruzada entre ambas 
representaciones. Comparado con t-SNE, \acs{UMAP} preserva 
mejor la estructura global de los datos, escala a 
conjuntos de datos significativamente más grandes y 
no tiene restricciones computacionales sobre la 
dimensión del espacio 
proyectado~\cite{mcinnes2018umap}.

Para que la reducción dimensional sea aplicable a 
datos textuales, los documentos o fragmentos de texto 
deben primero transformarse en vectores numéricos de 
alta dimensión. Los modelos de 
\textit{embeddings} textuales, como Sentence 
Transformers~\cite{reimers2019sbert}, proyectan 
oraciones o párrafos a un espacio vectorial denso en 
el que la distancia entre vectores refleja la 
similitud semántica entre los textos que representan. 
La combinación de \textit{embeddings} textuales con 
reducción dimensional mediante \acs{UMAP} permite visualizar 
grandes conjuntos de documentos como nubes de puntos 
en las que la proximidad espacial corresponde a la 
similitud de contenido.




\section{Procesamiento de Lenguaje Natural con \texorpdfstring{\acs{LLM}}{LLMs}}
La extracción de datos en documentos legales requería programar reglas manuales o entrenar modelos con miles de ejemplos específicos \cite{zhong-etal-2020-nlp}. Esta técnica resulta inviable para procesar expedientes de control gubernamental, cuyas estructuras y narrativas varían constantemente. Frente a este obstáculo, la presente investigación adopta los Grandes Modelos de Lenguaje como motor de extracción, aprovechando su capacidad para procesar semántica compleja sin requerir un reentrenamiento masivo por cada nuevo tipo de documento.


\subsection{Modelos y Extracción de Información}
Los \acs{LLM} poseen la capacidad de realizar tareas de extracción semántica avanzada sin requerir un reentrenamiento masivo. Mediante la técnica de aprendizaje con pocos ejemplos (\textit{few-shot learning})\cite{brown2020}. 

Para asegurar que la extracción de hitos, fechas y presupuestos mantenga una secuencia lógica, se emplea la estrategia de \textit{Prompting} denominada Cadena de Pensamiento (\textit{Chain-of-Thought}). Esta técnica induce al \acs{LLM} a descomponer un problema de análisis complejo en pasos de razonamiento intermedios antes de emitir una respuesta final, mejorando significativamente la fidelidad de la información extraída a partir de textos densos \cite{wei2022}.

\begin{figure}[H]
\centering
\includegraphics[width=0.8\textwidth]{Tesis_CS_202203/chain_of_thought.png}
\caption{Comparación entre el prompting estándar y el razonamiento por Cadena de Pensamiento, el cual muestra con claridad los pasos intermedios. Imagen extraída de \cite{wei2022}.}
\label{fig:cot_prompting}
\end{figure}



\subsection{Alucinaciones y Trazabilidad}
A pesar de su capacidad de razonamiento, la implementación de \acs{LLM} en entornos de alta responsabilidad enfrenta un riesgo crítico: las alucinaciones. \cite{ji2023} definen este fenómeno como la generación de contenido irreal que aparenta ser verdadero, clasificándolo en dos tipos: alucinaciones intrínsecas (cuando el texto generado contradice el documento fuente) y extrínsecas (cuando se genera información que no puede ser verificada en la fuente original). En una auditoría, alucinar la fecha , presupuesto o el lugar invalida por completo el análisis.

Frente al riesgo de las alucinaciones, la extracción de información no puede operar como una caja negra. \cite{hutchinson2024llm} advierten que para confiar en un \acs{LLM} dentro de un entorno de \textit{Visual Analytics}, sus respuestas deben estar ancladas a los datos originales. En el trabajo diario de un auditor, esto se traduce en una necesidad: el sistema debe crear un enlace visual directo entre el dato extraído por la máquina y el párrafo exacto del expediente. De esta forma, el analista no asume automáticamente que el modelo tiene la razón, sino que utiliza la interfaz para auditar la fuente, corregir errores y validar que la información sea real.

\subsection{Estructuración de Narrativas Complejas}
El objetivo de extraer entidades, tiempos y atributos no es generar bases de datos aisladas, sino reconstruir una narrativa compleja. En el contexto de esta investigación, una narrativa compleja se define como la progresión cronológica y causal de eventos burocráticos distribuidos a través de múltiples documentos (informes, adendas, resoluciones). 

Para modelar computacionalmente esta narrativa, el sistema extrae eventos. Un evento se define como una tupla estructurada que responde al qué, quién y cuándo de una acción documentada \cite{navas2022whenthefact, yao2022leven}. Al transformar el texto fragmentado en una secuencia de eventos, el \textit{Visual Analytics} logra mapear el flujo de la información. De esta forma, la interfaz revela visualmente cómo evolucionan los procesos y dónde se concentran las demoras, eliminando la necesidad de leer miles de documentos de forma secuencial.


\section{Consideraciones Finales}
El marco teórico desarrollado en este capítulo establece los pilares conceptuales para abordar el análisis de documentos gubernamentales. En primer lugar, se ha definido que el \textit{Visual Analytics} provee la arquitectura metodológica necesaria para mantener al analista humano en el centro del proceso de verificación. Paralelamente, las técnicas de visualización espaciotemporal y multidimensional proporcionan los mecanismos para representar la estructura cronológica y las métricas de las auditorías sin generar saturación visual.

Por su parte, se estableció que los Grandes Modelos de Lenguaje actúan como el motor de extracción capaz de procesar la semántica del texto y estructurar las narrativas complejas. Sin embargo, su implementación en escenarios de alta responsabilidad exige mecanismos de trazabilidad visual para mitigar el riesgo de alucinaciones y garantizar la fidelidad de la información.

En el siguiente capítulo se presentara el Estado del Arte. Se revisarán los trabajos enfocados en sistemas de extracción legal, herramientas visuales para texto y aplicaciones de los \acs{LLM}..
 %Inserta el capítulo 2
\chapter{Estado del Arte}
\label{cap:estado_arte}

\section{Consideraciones Iniciales}

En el presente capítulo se realiza una revisión exhaustiva de los trabajos relacionados con el análisis computacional, extracción de información y exploración visual de narrativas complejas en el contexto gubernamental y legal. Con el objetivo de establecer una base sólida de conocimiento, se han explorado diversas investigaciones que han abordado el problema desde diferentes perspectivas. Este capítulo se divide en tres enfoques principales: primero, los sistemas orientados a la extracción de información en documentos legales y gubernamentales; segundo, las herramientas de \textit{Visual Analytics} diseñadas para el análisis de texto y narrativas complejas temporales; y finalmente, las investigaciones centradas en el uso de Grandes Modelos de Lenguaje (\acs{LLM}) para el análisis documental y la mitigación de alucinaciones.

\section{Sistemas de Extracción de Información en Documentos Gubernamentales y Legales}

En esta sección describimos los trabajos que han aplicado técnicas de extracción de información, la gran parte de estas investigaciones intentan convertir documentos legales y gubernamentales en datos organizados.





En \cite{zhong-etal-2020-nlp}, Zhong et al. presentan un estudio detallado sobre la Inteligencia Artificial Legal (LegalAI), enfocándose en cómo el procesamiento de lenguaje natural puede ayudar a analizar sentencias y contratos. El trabajo propone un enfoque híbrido que mezcla modelos de lenguaje como \ac{BERT}, adaptado específicamente al dominio legal con métodos estadísticos y reglas lógicas para ganar interpretabilidad. Un problema técnico que identifican es la limitación de los 512 tokens en modelos como \acs{BERT}, esto es crítico porque un expediente legal real puede superar las 50,000 palabras, lo que causa que el modelo pierda el contexto global y falle en tareas de predicción, como ocurrió en el benchmark C-LJP. Además, los autores señalan que estos modelos suelen funcionar como cajas negras, lo que genera problemas éticos y falta de confianza en los resultados. Aunque mencionan el uso interno de grafos para relacionar crímenes y penas, el sistema no ofrece ninguna interfaz de \textit{Visual Analytics} para el usuario. Esta carencia es un punto clave para nuestra investigación, ya que refuerza la necesidad de herramientas que no solo procesen el texto, sino que permitan explorar visualmente narrativas largas y razonamientos complejos que los modelos actuales, por su propia arquitectura, todavía no pueden resolver de forma transparente.








En \cite{navas2022whenthefact} se presenta \textit{WhenTheFact}, una herramienta que integra la extracción legal con \textit{Visual Analytics} para identificar eventos y fechas en sentencias judiciales europeas. A diferencia de las propuestas basadas en modelos masivos, este sistema utiliza un pipeline de \ac{NLP} tradicional que combina \textit{CoreNLP} para el análisis sintáctico y \textit{FrameNet} para capturar la semántica de los verbos legales. El principal aporte visual es una interfaz interactiva en HTML que genera una línea de tiempo automática, permitiendo que el usuario navegue cronológicamente por los hechos del caso de forma secuencial. Sin embargo, el sistema tiene limitaciones importantes en la resolución de correferencias, lo que causa que muchas veces se extraigan pronombres en lugar de los nombres reales de los actores involucrados. Además, los autores señalan una fuerte dependencia de la estructura del documento original, si el formato del \ac{PDF} varía o es muy antiguo, la calidad de la extracción cae drásticamente. Este trabajo es un antecedente para nuestra investigación, ya que valida la utilidad de las líneas de tiempo para explorar narrativas, pero también expone cómo los métodos clásicos se quedan cortos ante problemas de coherencia y flexibilidad que los \acs{LLM} pueden resolver de forma más eficiente.












En \cite{yao2022leven} se presenta \textit{LEVEN}, un dataset diseñado para la detección de eventos legales en juicios penales. La metodología que proponen utiliza modelos de lenguaje como \acs{BERT} y DMBERT, apoyados en herramientas de segmentación como JIEBA y el uso de \ac{SBERT} para calcular la similitud semántica entre 108 categorías de eventos predefinidos. Un punto crítico que mencionan los autores es el problema de \textit{long-tail}, donde el rendimiento del modelo (F1) cae drásticamente en eventos con pocas muestras de entrenamiento. Además, el sistema tiene dificultades para integrar el contexto de oraciones cruzadas,como por ejemplo, no logra distinguir si el evento llamar es un reporte policial o una comunicación privada si no identifica correctamente la entidad a la que se contacta. Aunque el trabajo muestra cómo encadenar eventos (como un accidente que deriva en fuga y muerte) permite reclasificar un delito de tráfico a homicidio, el enfoque es algorítmico y carece de una interfaz de exploración para el usuario. Esta limitación es clave para mi investigación, ya que LEVEN demuestra que la extracción automática por sí sola todavía genera muchos falsos positivos y confusiones semánticas. Nuestra propuesta de \textit{Visual Analytics} busca llenar ese vacío, permitiendo que un analista valide y navegue estas narrativas complejas que los modelos actuales todavía no procesan con total precisión.










En \cite{shen2020hierarchical} se aborda el problema de la extracción jerárquica de eventos en documentos judiciales, proponiendo un modelo que busca estructurar automáticamente delitos y sentencias. La arquitectura utiliza \acs{BERT} para generar las representaciones vectoriales iniciales, pero añade un mecanismo de atención denominado \textit{pedal attention}. Este componente es clave porque permite conectar palabras que están semánticamente vinculadas pero muy alejadas en el texto, superando las limitaciones de las redes de grafos tradicionales que suelen fallar ante dependencias de larga distancia. Además, implementan un sistema de inferencia conjunta basado en matrices de probabilidad para asegurar que la estructura final del evento sea coherente y no se pierda la relación de causalidad entre un acto criminal y su resultado. El sistema no presenta una interfaz interactiva, pero utiliza mapas de calor (\textit{heatmaps}) para demostrar cómo el modelo distribuye la atención entre términos distantes. Esto resulta valioso para la interpretabilidad del sistema, permitiendo observar, por ejemplo, cómo se logra asociar correctamente una entidad con una acción específica en oraciones complejas. Sin embargo, el sistema aún presenta dificultades con la resolución de anáforas y ambigüedades cuando aparecen múltiples atributos en una sola frase. Este trabajo es un antecedente para nuestra investigación, ya que propone una forma robusta de estructurar eventos, pero refuerza la oportunidad de integrar estas capacidades en un sistema de \textit{Visual Analytics} que permita al usuario validar y explorar estas jerarquías de forma interactiva.












En \cite{feng2022legal} se presenta \textit{EPM}, un modelo diseñado para la predicción de juicios legales que busca determinar el delito y la condena a partir de los hechos de un caso. Su arquitectura utiliza Legal-BERT y una capa \ac{CRF} para extraer eventos, pero lo que realmente lo diferencia es el uso de restricciones lógicas rígidas dentro de la función de pérdida. Esto evita que el sistema cometa errores absurdos, como asignar un cargo que no tiene coherencia legal con el artículo de ley detectado. Además, el modelo soluciona el problema del ruido textual, logrando que el sistema se enfoque en la acción principal e ignore palabras secundarias que suelen confundir a modelos más simples. A pesar de su efectividad técnica, EPM funciona como un modelo de backend sin ninguna capa de \textit{Visual Analytics}. Los autores demuestran que extraer eventos mejora significativamente la precisión del juicio, pero los resultados se entregan como etiquetas de texto sin una interfaz que permita entender el contexto. Este punto para nuestra investigación, valida que la extracción de eventos es fundamental para analizar documentos judiciales, pero al mismo tiempo se necesita un sistema de \textit{Visual Analytics} que permita al analista humano explorar y validar estas narrativas complejas en lugar de solo recibir una predicción automática.








En \cite{xue2024leec} se introduce \textit{LEEC}, un dataset diseñado para la extracción de factores legales y extralegales (como datos demográficos) en documentos judiciales chinos. Para esta tarea, los autores evalúan tanto modelos tradicionales de extracción de eventos, como Doc2EDAG, frente a \acs{LLM} generales y especializados en el dominio legal. Un hallazgo crítico es la dificultad de los modelos para manejar documentos extensos debido a la limitación en la ventana de contexto, lo que provoca truncamientos y la pérdida de información clave. Además, el estudio reporta alucinaciones de formato donde los modelos fallan al generar salidas estructuradas en \ac{JSON} o confunden conceptos jurídicos básicos, por ejemplo, en el caso de un accidente de tránsito, GPT-3.5 no logró distinguir entre una compensación civil voluntaria y una multa penal. A pesar de la utilidad del dataset para estructurar información, la investigación no propone ningún componente de \textit{Visual Analytics}, limitándose a representaciones estáticas como gráficas de bosque para el análisis estadístico. Esta falta de una interfaz interactiva es un punto importante para nuestra investigación, ya que los errores detectados en la extracción de LEEC demuestran que no se puede confiar ciegamente en la automatización. Nuestra propuesta busca añadir esa capa visual que permita al usuario validar y navegar las narrativas complejas, mitigando los riesgos de alucinación o confusión semántica que los \acs{LLM} presentan.










En \cite{joshi2023ucreat} se propone \textit{U-CREAT}, un sistema no supervisado para la recuperación de casos legales basado en la extracción de eventos. La metodología utiliza un \textit{pipeline} de procesamiento de lenguaje natural que emplea \textit{dependency parsing} mediante la librería \textit{spaCy} para transformar oraciones en grafos dirigidos, extrayendo tuplas de eventos compuestas por sujetos, verbos y objetos. Los autores descartan el uso de arquitecturas basadas en \textit{Transformers} como \acs{BERT} para la extracción final, debido a las limitaciones en el \textit{input token length} y el elevado tiempo de inferencia que penaliza el rendimiento en documentos legales extensos. En su lugar, el sistema utiliza modelos estadísticos como \ac{BM25} y similitud de Jaccard para medir la relevancia entre casos. A pesar de extraer con precisión acciones como \texttt{(statement, forward, police)}, el sistema falla al no capturar la naturaleza secuencial de la narrativa legal, al tratar los eventos de forma aislada, se pierde el flujo temporal y el orden cronológico de los hechos. Además, U-CREAT carece de una interfaz de \textit{Visual Analytics}, limitándose a gráficas estáticas de rendimiento (\textit{F1 Score}) y tiempo de inferencia. Esta falta de una capa de exploración visual interactiva es clave para nuestra investigación, ya que la incapacidad del modelo para reconstruir la secuencia de los hechos justifica nuestra propuesta que permita al analista validar y navegar la estructura espacio temporal de estas narrativas complejas.

En \cite{pereira2023inacia} se presenta \textit{INACIA}, un sistema diseñado para automatizar el análisis de expedientes en tribunales de Brasil. A diferencia de un sistema de \textit{Visual Analytics}, este se enfoca en la parte operativa: procesar denuncias en \ac{PDF} y generar propuestas legales. Su arquitectura es un modelo híbrido que utiliza GPT-4 (con una ventana de 32k tokens) y un flujo \ac{RAG} que se apoya en motores de búsqueda como \acs{BM25} y el reranker \textit{NeuralSearchX}. Un punto interesante de su metodología es que para decisiones críticas, como el \textit{Periculum in mora}, no confían ciegamente en el modelo de lenguaje, sino que usan reglas fijas (\textit{hard-coded}) para validar datos exactos antes de redactar cualquier justificación. Sin embargo, el sistema todavía tiene fallas importantes. Los autores mencionan que su \textit{recall} es de apenas 40.1\%, lo que significa que el modelo omite mucha información que un humano sí consideraría. También se nota que el sistema se confunde al manejar documentos muy largos con leyes contradictorias, lo que genera alucinaciones. Un ejemplo de esto fue su desempeño con las licitaciones de IBAMA: aunque extrajo datos simples con un 100\% de precisión, falló en los razonamientos legales más complejos. Finalmente, es importante resaltar que INACIA no tiene una interfaz visual; todo se maneja por código o \acp{API}. Esta falta de una herramienta de exploración visual es justamente lo que justifica nuestra propuesta que busca que el usuario pueda navegar e interactuar con estas narrativas legales y no solo recibir el texto procesado.





Si bien todos estos enfoques y sistemas demuestran una alta viabilidad técnica para extraer eventos, hitos y entidades en documentos oficiales, presentan una limitación compartida respecto a las necesidades de las auditorías gubernamentales: se centran exclusivamente en la automatización textual y algorítmica. Al carecer de un componente interactivo de \textit{Visual Analytics}, no permiten al analista validar de forma visual y espacial el origen de cada dato extraído, un requisito indispensable para generar confianza y trazabilidad.









\section{Herramientas de \textit{Visual Analytics} para Texto y Narrativas Complejas}

En esta sección exploramos las herramientas de \textit{Visual Analytics} que han sido diseñadas para el análisis de texto; la mayoría de estos trabajos buscan que el usuario no solo reciba datos extraídos, sino que pueda interactuar con la información. El objetivo de estos autores es permitir que los analistas descubran por sí mismos patrones o cambios en el tiempo, integrando técnicas de Procesamiento de Lenguaje Natural con interfaces visuales que facilitan la exploración de grandes volúmenes de documentos






En \cite{knittel2021} presentan \textit{PyramidTags}, un sistema de \textit{Visual Analytics} para explorar grandes colecciones de noticias sin conocer el contenido de antemano. Este trabajo no usa \acs{LLM}, sino que se basa en un flujo estadístico que emplea \textit{tf-idf} para extraer etiquetas y el algoritmo de optimización por enjambre de partículas (\ac{PSO}) para acomodarlas en un mapa 2D. El posicionamiento se define mediante una función de energía que busca equilibrar el tiempo, la relevancia y evitar que las palabras se superpongan. Visualmente, el sistema usa una pirámide interactiva donde el eje horizontal indica cuándo ocurrió el evento y el vertical cuánto duró. Si bien la herramienta permite reconstruir historias reales, como que ocurrió con los reportes de la cumbre de Singapur en 2018, los autores identifican fallos importantes como los falsos amigos. Esto pasa porque el algoritmo agrupa términos que no tienen relación semántica solo porque aparecieron en fechas cercanas, como juntar noticias de funerales con escándalos mediáticos. Además, el diseño deja mucho espacio vacío y las etiquetas se solapan al interactuar con ellas. Estos problemas demuestran que los métodos estadísticos tienen un límite al no entender el contexto, lo que justifica el uso de \acs{LLM} en nuestra tesis para lograr una extracción de narrativas mucho más precisa.





En \cite{krause2023} proponen un sistema de \textit{Visual Analytics} diseñado para identificar puntos de cambio en series de tiempo de grandes colecciones de texto. La arquitectura se basa en un modelo de \textit{Rolling \ac{LDA}} para extraer temas dinámicos y utiliza la similitud del coseno para detectar matemáticamente cuándo ocurre una disrupción en el vocabulario. El sistema ofrece un tablero con nubes de palabras adaptativas que incluyen \textit{scarfplots} y gráficos de barras tipo mariposa (\textit{butterfly charts}) para detallar qué términos provocaron el cambio estadístico. Si bien la herramienta fue eficaz detectando eventos abruptos, como la escasez de oxígeno en India o escándalos políticos puntuales durante la pandemia, los autores admiten fallos en el rastreo de conceptos dinámicos. El algoritmo, al ser estrictamente estadístico falla cuando una palabra clave aumenta su frecuencia de manera constante, ya que la métrica de similitud no está diseñada para procesar crecimientos prolongados. Esta limitación de los modelos tradicionales para entender progresiones complejas refuerza nuestra propuesta, donde el uso de Grandes Modelos de Lenguaje (\acs{LLM}) permite capturar la evolución narrativa y no solo matemática.










En \cite{chen2019supporting} proponen un sistema de \textit{Visual Analytics} para la síntesis de historias basado en datos de redes sociales y componentes espacio temporales. Para el procesamiento no utiliza modelos de lenguaje, sino usa técnicas estadísticas como \textit{Latent Dirichlet Allocation} (\acs{LDA}) y algoritmos de reducción de dimensiones como t-SNE y \ac{MDS} para organizar los temas. El sistema ofrece una interfaz de vistas múltiples coordinadas que incluye flujos de palabras clave (\textit{keyword flow}), mapas geográficos y diseños anidados (\textit{nested layouts}) para navegar por la información en diferentes niveles. Sin embargo, los autores identifican fallos, por ejemplo, el modelo \acs{LDA} falla al procesar textos muy breves, como los tuits, porque no encuentra suficientes palabras significativas por documento. Esto los obligó a usar un parche en su investigación sobre el Brexit, teniendo que agrupar mensajes manualmente en mega documentos antes de poder analizarlos. Además, el sistema tiene dificultades para evitar la redundancia en los reportes y no logra mitigar los sesgos del analista que termina ignorando partes de la base de datos. Este trabajo nos sirve como antecedente porque valida la importancia de las interfaces interactivas para explorar historias, pero también que los métodos estadísticos tradicionales son muy limitados ante el lenguaje ruidoso. Esto justifica nuestra propuesta de usar \acs{LLM}, los cuales permiten capturar la semántica de forma más profunda y sin depender de agrupar manualmente los documentos.








En \cite{rajabiyazdi2024textvista} proponen \textit{TextVista}, una herramienta de \textit{Visual Analytics}  diseñada para explorar patrones, sentimientos y entidades en grandes volúmenes de texto no estructurado con referencias temporales. La arquitectura del sistema es híbrida, utiliza modelos \textit{Transformers} como \acs{SBERT} para generar \textit{embeddings} semánticos y \acs{BERT} para la clasificación de sentimientos, además de una red Bi-LSTM-\acs{CRF} para el reconocimiento de entidades. Para la exploración visual, el sistema emplea una interfaz de vistas múltiples coordinadas que incluye un panel de frecuencias temporales, un dial de entidades radial y un gráfico de dispersión 2D generado mediante \ac{PCA} para visualizar el espacio semántico. Si bien la herramienta demostró su utilidad en pruebas reales permitiendo a los analistas identificar redes de \textit{bots} en temas de criptomonedas y desinformación, los autores reportan una falla crítica en la explicabilidad del modelo. Los usuarios manifestaron frustración al no comprender el razonamiento detrás del agrupamiento semántico, señalando que el sistema funciona como una caja negra que muestra la separación matemática de los datos pero no su lógica narrativa. Este antecedente es importante para nuestro trabajo, ya que valida el uso de arquitecturas basadas en \textit{Transformers} para el procesamiento de texto, pero también resalta la necesidad de mejorar la transparencia de los datos. Nuestra propuesta busca cubrir este vacío utilizando las capacidades de razonamiento de los \acs{LLM} para facilitar la interpretación de las narrativas complejas.









En \cite{zhao2023contextwing} presentan \textit{ContextWing}, un sistema de \textit{Visual Analytics} diseñado para la comparación binaria y evolutiva de patrones secuenciales en flujos de datos de redes sociales. La metodología del sistema es híbrida, utiliza modelos basados en \textit{Transformers} como \acs{BERT} para el análisis de sentimiento y técnicas dinámicas como \textit{BERTopic}, combinando \ac{UMAP} y K-Means para asegurar la coherencia temporal de los temas en tiempo real. El componente visual es su metáfora de alas bilaterales, donde una palabra clave actúa como eje y el contexto anterior y posterior se despliega en capas simétricas para comparar la percepción entre dos entidades. Si bien la herramienta demostró su eficacia durante el debate presidencial de EE. UU. de 2020, detectando en minutos el cambio de narrativa desde temas económicos hacia discusiones sobre inmigración, los autores reconocen una limitación técnica importante: el diseño está restringido a comparaciones de pares (\textit{pair-wise}). Escalar este modelo a múltiples flujos de datos es un problema que volvería redundante la estructura visual. Además, la eficacia del sistema depende de que los datos sean heterogéneos, perdiendo utilidad en contextos donde los datos son muy similares entre sí. Esta limitación de comparar solo un par a la vez nos deja un vacío que queremos llenar. Nuestra propuesta usa \acs{LLM} para explorar varias historias y puntos de vista al mismo tiempo, dándonos mucha más libertad visual de la que permiten los gráficos tradicionales.







En \cite{lacava2022} se presenta \textit{LawNet-Viz}, un sistema híbrido que integra el análisis de redes y modelos de lenguaje para la exploración de documentos legales. Su arquitectura emplea reglas sintácticas para identificar citas entre artículos y utiliza el modelo \textit{LamBERTa}, es un \acs{BERT} especializado en el dominio legal italiano para el cálculo de similitudes semánticas. El componente visual principal consiste en un grafo interactivo donde los artículos se representan como nodos y las aristas como referencias explícitcas, aplicando métricas de centralidad como \textit{PageRank} y \textit{Betweenness} para resaltar la relevancia de cada norma. La interfaz permite además una navegación por vecindarios de hasta dos saltos (\textit{2-hop}) para inspeccionar el contexto legal de cada artículo. Aunque el sistema logró mapear con éxito el Código Civil Italiano, los autores identifican limitaciones en la generalización del extractor, el cual depende de configuraciones manuales específicas y descarta las referencias externas al documento analizado. Asimismo, se reportan dificultades de \textit{few-shot learning} debido a la alta cantidad de clases que el modelo debe manejar. Este trabajo es un antecedente relevante porque valida la utilidad de los grafos en el dominio legal, pero evidencia que la dependencia de reglas rígidas limita la flexibilidad del análisis. Esto justifica nuestra propuesta de utilizar \acs{LLM} para una exploración de narrativas más contextual que no dependa de estructuras de extracción predefinidas para cada dominio.




En \cite{resck_legalvis} proponen \textit{LegalVis}, un sistema de \textit{Visual Analytics} diseñado para explorar precedentes en documentos legales. La metodología es híbrida, utiliza reglas de extracción y un modelo de clasificación basado en \textit{\acs{TFIDF}} junto con una Máquina de Vectores de Soporte (\ac{SVM}). Durante las pruebas, los autores notaron que modelos más complejos como \textit{Longformer} rendían peor debido a la falta de preentrenamiento en el lenguaje legal específico del portugués. Visualmente, la herramienta usa un enfoque de vistas vinculadas que incluye gráficos temporales, mapas de similitud y un lector de documentos que usa el método LIME para dar interpretabilidad al modelo. Aunque el sistema logró identificar citas que las reglas clásicas ignoraban, tiene limitaciones de escala. Debido al desbalance de los datos, solo pudo procesar una pequeña parte de los precedentes disponibles y mostró fallas de precisión en contextos poco relacionados. Este trabajo es un antecedente para nuestra investigación, ya que confirma que el éxito de los \acs{LLM} en el área legal depende mucho de su especialización. Además, valida la necesidad de sistemas que superen los problemas de escala y desbalance mediante un razonamiento contextual más robusto.







En \cite{yeh2025storyribbons} proponen \textit{Story Ribbons}, un sistema de \textit{Visual Analytics} que utiliza \acs{LLM} para estructurar y visualizar narrativas en textos literarios extensos. La arquitectura emplea un enfoque multimodelo, utiliza \textit{gpt-4o-mini} para la extracción de datos y \textit{claude-3-5-sonnet} para ejecutar bucles de corrección que consolidan elementos duplicados. Para mitigar las alucinaciones en las citas textuales, los autores implementaron una regla de validación de cadenas (\textit{exact string match}) que descarta cualquier texto generado que no coincida con la fuente original. El componente visual principal consiste en diagramas de cintas (\textit{ribbon plots}) donde el eje vertical es personalizable para rastrear sentimientos, ubicaciones o niveles de importancia de los personajes a lo largo del tiempo. Sin embargo, el estudio identifica limitaciones relacionadas con el manejo de contextos largos y la resolución de entidades, por ejemplo, el sistema falló al no reconocer que Polifemo y el Cíclope eran el mismo personaje en la Odisea. Además, se reportó que las síntesis literarias tienden a ser superficiales, perdiendo los matices psicológicos de los personajes. Esta investigación es un antecedente importante para nuestra investigación, ya que demuestra la potencia de los \acs{LLM} para estructurar historias. Esto refuerza nuestra propuesta de desarrollar un sistema de \textit{Visual Analytics} que logre una exploración más profunda y coherente mediante el uso de \acs{LLM} con mayor capacidad de razonamiento contextual.













\section{Uso de \texorpdfstring{\acs{LLM}}{LLMs} para Análisis Documental e Interacción de Datos}

La irrupción de los Grandes Modelos de Lenguaje (\acs{LLM}) ha movido el foco del análisis documental desde tareas aisladas de extracción hacia flujos completos de comprensión, estructuración e interacción. En este eje, el aporte ya no es solo ``extraer mejor'', sino integrar razonamiento, consulta y validación en una misma cadena de trabajo.

En \cite{ke2025document} presentan una revisión amplia sobre inteligencia documental basada en \acs{LLM}. Los autores explican que hoy el campo se está moviendo en tres líneas principales: uso de \textit{Retrieval-Augmented Generation} para responder con evidencia, manejo de contexto largo para trabajar con documentos extensos y ajuste fino para dominios específicos. El aporte más útil para esta tesis es que el problema ya no se entiende como ``extraer un campo con buena precisión'', sino como mantener coherencia, fidelidad y trazabilidad en todo el flujo de trabajo. También señalan una limitación importante: cuando el documento crece en tamaño, aumenta la dificultad para conservar contexto y verificar cada salida del modelo. Esto conecta directamente con nuestra propuesta, porque muestra que no basta con generar respuestas; también se necesita una capa de revisión para que el analista pueda comprobar de dónde sale cada resultado.

En \cite{yue2023event} proponen un marco cooperativo con \acs{LLM} para generar vistas judiciales a partir de eventos extraídos en casos penales. La metodología divide el proceso en etapas: primero se identifican eventos clave del expediente (actores, acciones y relaciones), y luego esos eventos se usan para construir una narrativa judicial más ordenada y menos ambigua que el texto libre. El aporte es útil porque obliga al modelo a trabajar sobre una estructura intermedia y no solo sobre párrafos sueltos, lo que mejora la coherencia del resultado final. El problema es que la calidad de la salida depende mucho de la extracción inicial: si un evento se interpreta mal, ese error se arrastra a toda la narrativa. Un ejemplo que discuten es que, cuando se pierde la secuencia correcta entre hechos vinculados, la interpretación del caso puede cambiar de forma importante. Además, aunque el enfoque organiza mejor la información legal, sigue siendo principalmente textual y no ofrece una interfaz visual para auditar cada paso, que es justamente el vacío que nuestra propuesta de \acs{VA} busca cubrir.

En \cite{wei2022} muestran que la técnica \textit{Chain-of-Thought} mejora el rendimiento de los \acs{LLM} cuando la tarea exige razonamiento paso a paso. La metodología consiste en guiar al modelo con ejemplos donde no solo se da la respuesta final, sino también el proceso intermedio que lleva a esa respuesta. En sus experimentos, reportan mejoras claras en tareas como GSM8K usando modelos grandes como PaLM 540B, lo que confirma que explicitar pasos ayuda a reducir errores en problemas complejos. Para el análisis documental, este aporte es valioso porque permite separar mejor la fase de interpretación de la fase de salida estructurada. Sin embargo, los autores también dejan ver una limitación: razonar en varios pasos no garantiza por sí solo que el contenido sea fiel a la fuente, ya que el modelo puede construir una cadena lógica convincente pero basada en premisas incorrectas. Por eso, en nuestro contexto, \textit{Chain-of-Thought} es útil como estrategia de procesamiento, pero necesita complementarse con mecanismos de verificación y revisión visual para validar cada conclusión antes de usarla en un entorno institucional.

En \cite{zhong2017seq2sql} proponen \textit{Seq2SQL}, un modelo para traducir preguntas en lenguaje natural a consultas SQL ejecutables. La metodología combina una red neuronal con aprendizaje por refuerzo, donde el sistema no solo aprende la forma de la consulta, sino que también se optimiza usando el resultado de ejecución para reducir errores semánticos. En sus pruebas sobre WikiSQL, los autores muestran una mejora importante frente a enfoques anteriores, especialmente en exactitud de ejecución, lo que demuestra que validar contra la base de datos real ayuda a generar consultas más correctas. Este trabajo es relevante para nuestra tesis porque conecta directamente con la interacción entre lenguaje natural y datos estructurados. Aun así, también presenta una limitación: funciona mejor en consultas de complejidad moderada y en esquemas relativamente controlados, por lo que al crecer la complejidad del dominio aparecen más ambigüedades y fallos en la traducción. En otras palabras, \textit{Seq2SQL} resuelve bien la conversión inicial a SQL, pero todavía requiere apoyo adicional cuando se busca trazabilidad completa y validación en escenarios institucionales más exigentes.

En \cite{ye2024dataframe} desarrollan un framework que usa \acs{LLM} para generar consultas Pandas sobre dataframes, con énfasis en privacidad: el modelo recibe solo los metadatos y nombres de columnas, no los datos reales. La metodología es simple pero efectiva: con esa información limitada, el modelo genera código Python que el usuario ejecuta localmente, evitando la exposición de información sensible. En sus evaluaciones, logran porcentajes altos de exactitud en datasets públicos como WikiSQL y UCI-DataFrameQA. El valor para esta tesis es que demuestra cómo permitir interacción con datos sin comprometer privacidad, algo crítico en contextos institucionales como la Contraloría. Sin embargo, el enfoque tiene una debilidad clara: si la estructura del dataframe es compleja o ambigua, el modelo también genera consultas ambiguas o incorrectas, porque no puede ver los datos para validar su interpretación. Un ejemplo que mencionan es que columnas con nombres crípticos o contextos no obvios frecuentemente producen consultas equivocadas. Por eso, nuestra propuesta de \acs{VA} se complementa bien con este tipo de herramientas: el usuario ejecuta una consulta generada, ve los resultados en el gráfico y puede revisar si la lógica fue correcta antes de usar esos datos en un reporte oficial.

En \cite{dibia2023lida} presentan LIDA, un sistema que usa \acs{LLM} para automatizar la creación de visualizaciones a partir de datos tabulares. La metodología trabaja por etapas: primero resume el dataset, luego propone objetivos de análisis y finalmente genera o refina el código de la visualización según la intención del usuario. Este enfoque es importante porque reduce la barrera técnica para construir gráficos y acelera la exploración inicial de datos. Para nuestra tesis, su aporte principal es mostrar que un \acs{LLM} puede actuar como puente entre una pregunta en lenguaje natural y una salida visual concreta. Sin embargo, el propio enfoque deja una limitación clara: que el gráfico se vea bien no garantiza que la interpretación sea correcta ni que respete totalmente la fuente original. En la práctica, esto puede producir visualizaciones plausibles pero con decisiones de agregación o filtrado que el usuario no detecta a simple vista. Por eso, este tipo de sistemas necesita una capa de validación explícita, y ahí es donde nuestra propuesta de \acs{VA} aporta trazabilidad para revisar cómo se construyó cada resultado.

En \cite{ji2023} presentan una revisión amplia sobre alucinaciones en generación de lenguaje natural y proponen una distinción útil entre alucinaciones intrínsecas y extrínsecas. La metodología del trabajo es de tipo survey: comparan tareas, métricas y causas del problema en distintos escenarios de uso de modelos generativos. El aporte clave para esta tesis es que ayuda a entender que no todos los errores tienen el mismo impacto: en las intrínsecas, el modelo contradice directamente la fuente, y en las extrínsecas produce contenido no verificable aunque suene convincente. Esta diferencia es muy importante en análisis documental, porque en ambos casos se compromete la confiabilidad del resultado final. Los autores también muestran que el problema no se resuelve solo con modelos más grandes, ya que la alucinación depende también del contexto, del prompt y de la forma en que se valida la salida. En nuestro caso, este trabajo refuerza la necesidad de combinar extracción automática con mecanismos de revisión humana, para detectar a tiempo información inventada o débilmente sustentada antes de que llegue a una decisión institucional.

En \cite{hutchinson2024llm} analizan cómo integrar \acs{LLM} en flujos de \acs{VA} para apoyar tareas de exploración y toma de decisiones. La metodología del trabajo es de revisión y síntesis de oportunidades y riesgos: muestran dónde los modelos ayudan de verdad (por ejemplo, traducir preguntas en lenguaje natural a operaciones analíticas) y dónde todavía fallan (confiabilidad, explicabilidad y control del usuario). El aporte más valioso para esta tesis es que plantean al \acs{LLM} como asistente dentro del proceso analítico, no como reemplazo del analista. También remarcan una limitación práctica: cuando la interacción se vuelve muy automática, el usuario puede aceptar resultados plausibles sin revisar su sustento, lo que aumenta el riesgo de error en contextos sensibles. En ese sentido, el trabajo refuerza que la utilidad real aparece cuando la interfaz permite contrastar evidencia, iterar hipótesis y mantener trazabilidad de cada paso. Esa idea coincide con nuestra propuesta, porque buscamos que el analista conserve control sobre el proceso y pueda validar visualmente las salidas antes de convertirlas en insumo de auditoría.

En \cite{ha2024confides} presentan ConFides, un sistema de \acs{VA} para analizar resultados de reconocimiento automático del habla y revisar qué tan confiables son las transcripciones. La metodología combina las salidas del modelo con puntajes de confianza y las lleva a una interfaz con vistas coordinadas, donde el usuario puede ver rápidamente qué partes del texto son más inciertas y necesitan revisión. El aporte principal es que no oculta la incertidumbre, sino que la vuelve visible para apoyar una validación humana más informada. Esto es valioso para esta tesis porque sigue la misma lógica que buscamos: usar modelos automáticos, pero con control del analista en cada paso crítico. Una limitación del enfoque es que depende de la calidad de esos puntajes de confianza; si están mal calibrados, la interfaz también puede guiar mal la revisión. Además, está centrado en transcripción de audio, así que su adaptación a expedientes documentales requiere ajustes en el tipo de evidencia mostrada. Aun así, su propuesta confirma que hacer visible la incertidumbre mejora la trazabilidad y reduce el riesgo de aceptar resultados automáticos sin verificación.

\section{Consideraciones Finales}

La revisión presentada en este capítulo permite identificar tres tendencias y una brecha compartida que motivan la propuesta de esta tesis.

En la primera sección se analizaron sistemas de extracción de información en documentos legales y gubernamentales. Los enfoques basados en modelos como \acs{BERT} y técnicas de \ac{NLP} tradicional~\cite{zhong-etal-2020-nlp, navas2022whenthefact, yao2022leven, shen2020hierarchical, feng2022legal, joshi2023ucreat} logran extraer entidades y eventos con precisión aceptable, pero operan predominantemente a nivel de documento individual y dependen de configuraciones específicas por dominio. Los trabajos que incorporan \acs{LLM}~\cite{xue2024leec, pereira2023inacia} demuestran mayor flexibilidad, aunque reportan alucinaciones de formato y baja cobertura en razonamientos complejos. Ninguno de estos sistemas ofrece una interfaz de \textit{Visual Analytics} que permita al analista validar visualmente los resultados.

En la segunda sección se revisaron herramientas de \ac{VA} para texto y narrativas. Los sistemas basados en métodos estadísticos como \acs{TFIDF}, \ac{LDA} o \ac{PSO}~\cite{knittel2021, krause2023, chen2019supporting} logran visualizar patrones temporales, pero fallan al capturar la semántica profunda del texto. Las herramientas que adoptan modelos \textit{Transformers}~\cite{rajabiyazdi2024textvista, zhao2023contextwing, lacava2022, resck_legalvis} mejoran la representación semántica, aunque reportan problemas de explicabilidad y limitaciones de escala. Story~Ribbons~\cite{yeh2025storyribbons} es el único sistema que combina \acs{LLM} con \ac{VA} para narrativas, pero está orientado a textos literarios y no aborda documentos institucionales fragmentados ni la dimensión geoespacial.

En la tercera sección se examinó el uso de \acs{LLM} para análisis documental. Los trabajos revisados~\cite{ke2025document, yue2023event, wei2022, zhong2017seq2sql, ye2024dataframe, dibia2023lida} muestran que los \acs{LLM} permiten razonamiento encadenado, traducción de lenguaje natural a consultas estructuradas y generación automática de visualizaciones. Sin embargo, Ji~et~al.~\cite{ji2023} y Hutchinson~et~al.~\cite{hutchinson2024llm} advierten que la alucinación persiste como riesgo inherente y que la integración de \acs{LLM} en flujos de \ac{VA} requiere mecanismos explícitos de trazabilidad y control del analista. ConFides~\cite{ha2024confides} valida que hacer visible la incertidumbre del modelo mejora la confianza del usuario en el proceso de revisión.

La Tabla~\ref{tab:comparativa_ea} sintetiza la comparación entre los sistemas revisados y nuestra propuesta. Se observa que ningún trabajo previo reúne simultáneamente las seis capacidades que esta tesis integra: extracción mediante \acs{LLM}, interfaz de \ac{VA} con vistas coordinadas, soporte multi-documento, trazabilidad al fragmento fuente, mitigación de alucinaciones y representación de la dimensión espacio-temporal.

\begin{table}[htbp]
\centering
\caption{Comparación de los sistemas revisados. \checkmark{} = presente, $\circ$ = parcial, $\times$ = ausente. Columnas: \textbf{T}~=~Técnica (\textit{L}:~\acs{LLM}, \textit{B}:~\acs{BERT}/Transformers, \textit{E}:~Estadístico, \textit{R}:~Reglas/Tradicional, \textit{N}:~Red neuronal), \textbf{VA}~=~Visual Analytics, \textbf{MD}~=~Multi-documento, \textbf{Tz}~=~Trazabilidad, \textbf{MA}~=~Mitigación de alucinaciones, \textbf{ET}~=~Espacio-temporal.}
\label{tab:comparativa_ea}
\renewcommand{\arraystretch}{1.15}
\footnotesize
\begin{tabular}{@{}l l c c c c c c@{}}
\toprule
& \textbf{Trabajo} & \textbf{T} & \textbf{VA} & \textbf{MD} & \textbf{Tz} & \textbf{MA} & \textbf{ET} \\
\midrule
\multicolumn{8}{l}{\textit{Sistemas de extracción en documentos legales y gubernamentales}} \\
\midrule
1 & Zhong et al.~\cite{zhong-etal-2020-nlp}      & B & $\times$ & $\times$ & $\times$ & $\times$ & $\times$ \\
\rowcolor{rowgray}
2 & WhenTheFact~\cite{navas2022whenthefact}       & R & \checkmark & $\times$ & $\times$ & $\times$ & \checkmark \\
3 & LEVEN~\cite{yao2022leven}                     & B & $\times$ & $\times$ & $\times$ & $\times$ & $\times$ \\
\rowcolor{rowgray}
4 & Shen et al.~\cite{shen2020hierarchical}       & B & $\times$ & $\times$ & $\times$ & $\times$ & $\times$ \\
5 & EPM~\cite{feng2022legal}                      & B & $\times$ & $\times$ & $\times$ & $\times$ & $\times$ \\
\rowcolor{rowgray}
6 & LEEC~\cite{xue2024leec}                       & L & $\times$ & $\times$ & $\times$ & $\times$ & $\times$ \\
7 & U-CREAT~\cite{joshi2023ucreat}                & R & $\times$ & \checkmark & $\times$ & $\times$ & $\times$ \\
\rowcolor{rowgray}
8 & INACIA~\cite{pereira2023inacia}               & L & $\times$ & \checkmark & $\times$ & $\circ$  & $\times$ \\
\midrule
\multicolumn{8}{l}{\textit{Herramientas de Visual Analytics para texto y narrativas}} \\
\midrule
9  & PyramidTags~\cite{knittel2021}               & E & \checkmark & \checkmark & $\times$ & $\times$ & \checkmark \\
\rowcolor{rowgray}
10 & Krause et al.~\cite{krause2023}              & E & \checkmark & \checkmark & $\times$ & $\times$ & \checkmark \\
11 & Chen et al.~\cite{chen2019supporting}         & E & \checkmark & \checkmark & $\times$ & $\times$ & \checkmark \\
\rowcolor{rowgray}
12 & TextVista~\cite{rajabiyazdi2024textvista}     & B & \checkmark & \checkmark & $\times$ & $\times$ & \checkmark \\
13 & ContextWing~\cite{zhao2023contextwing}        & B & \checkmark & \checkmark & $\times$ & $\times$ & \checkmark \\
\rowcolor{rowgray}
14 & LawNet-Viz~\cite{lacava2022}                 & B & \checkmark & \checkmark & $\times$ & $\times$ & $\times$ \\
15 & LegalVis~\cite{resck_legalvis}               & E & \checkmark & \checkmark & $\times$ & $\times$ & \checkmark \\
\rowcolor{rowgray}
16 & Story Ribbons~\cite{yeh2025storyribbons}     & L & \checkmark & \checkmark & $\circ$  & $\circ$  & \checkmark \\
\midrule
\multicolumn{8}{l}{\textit{Uso de LLMs para análisis documental e interacción}} \\
\midrule
17 & Yue et al.~\cite{yue2023event}               & L & $\times$ & $\times$ & $\times$ & $\times$ & $\times$ \\
\rowcolor{rowgray}
18 & Seq2SQL~\cite{zhong2017seq2sql}               & N & $\times$ & $\times$ & $\times$ & $\times$ & $\times$ \\
19 & DataFrame QA~\cite{ye2024dataframe}           & L & $\times$ & $\times$ & $\times$ & $\times$ & $\times$ \\
\rowcolor{rowgray}
20 & LIDA~\cite{dibia2023lida}                     & L & \checkmark & $\times$ & $\times$ & $\times$ & $\times$ \\
21 & ConFides~\cite{ha2024confides}                & N & \checkmark & \checkmark & \checkmark & $\times$ & $\times$ \\
\midrule
\rowcolor{rowgray}
   & \textbf{Nuestra propuesta}                    & \textbf{L} & \checkmark & \checkmark & \checkmark & \checkmark & \checkmark \\
\bottomrule
\end{tabular}
\end{table}

La brecha que se observa en la tabla motiva directamente el diseño del sistema descrito en el siguiente capítulo. En él se detalla cómo la propuesta articula un pipeline de extracción basado en \acs{LLM} con corrección automática de inconsistencias, un mecanismo de trazabilidad que vincula cada dato con su fragmento de origen, y un conjunto de vistas interactivas coordinadas que integran las dimensiones temporal, geoespacial y narrativa de los procesos documentados.
        

 %Inserta el capítulo 3
\chapter{Propuesta} % aqui hay que nombrar al sistema :P
\label{cap:propuesta}

\section{Consideraciones Iniciales}

Este capítulo detalla el diseño y la implementación de nuestro 
sistema para ayudar a los analistas a explorar procesos complejos 
documentados en texto. Nuestro sistema busca resolver tres retos 
principales. Primero, la dificultad de reconstruir la narrativa 
de un caso a partir de documentos fragmentados. Segundo, la 
validación de los datos extraídos. Tercero, la integración de 
las dimensiones temporal, espacial y narrativa en una sola 
interfaz. Para ello, definimos primero los objetivos de diseño 
y las tareas analíticas que guían la propuesta. Luego, explicamos 
la implementación del sistema dividida en tres etapas clave, 
extracción, agregación y exploración interactiva.

A lo largo del capítulo usamos tres conceptos centrales. Un 
\textbf{caso} es la unidad de análisis del sistema, una 
organización, proyecto u obra sobre la que se dispone de 
múltiples documentos. Un \textbf{informe} es cada uno de esos 
documentos, que registra el estado del caso en un punto 
particular de su evolución. Una \textbf{variable} es cualquier 
atributo del caso que el sistema extrae y rastrea a lo largo 
de los informes, como el presupuesto ejecutado, los hallazgos 
identificados o el tipo de acción registrada.

\section{Objetivos y Tareas}

A partir de la revisión de la literatura y del análisis de los 
problemas identificados previamente, se definieron los siguientes 
objetivos de diseño y tareas analíticas que guían el desarrollo 
del sistema.

\subsection{Objetivos de Diseño}

\textbf{OD1. Analizar la evolución, ubicación y variables de 
los casos.} El análisis requiere examinar cómo cambian los casos 
a lo largo del tiempo, identificando sus ubicaciones geográficas 
en distintos niveles y comparando sus variables clave a partir 
de la información extraída de múltiples informes.

\textbf{OD2. Identificar similitudes y agrupaciones entre casos.}
El análisis requiere examinar el conjunto de casos para detectar 
cuáles comparten características comunes, e identificar qué tan 
atípico o recurrente es un caso dentro del conjunto.

\textbf{OD3. Verificar el origen de cada dato extraído.}
La confiabilidad del análisis requiere examinar que cada dato 
extraído proviene del documento original, permitiendo identificar 
posibles errores.

\subsection{Tareas Analíticas}

\textbf{T1. Reconstruir la evolución temporal de los casos.}
El sistema permitirá organizar cronológicamente los informes de 
cada caso para visualizar su desarrollo en el tiempo. \textbf{[OD1]}

\textbf{T2. Comparar el progreso entre casos.}
El sistema permitirá analizar la evolución de varios casos en 
paralelo, comparando sus informes y variables para identificar 
similitudes o diferencias en su desarrollo. \textbf{[OD1]}

\textbf{T3. Localizar los casos geográficamente.} 
El sistema permitirá visualizar la ubicación de los casos a 
nivel nacional, regional y local para examinar su concentración 
y detalle en diferentes niveles. \textbf{[OD1]}

\textbf{T4. Explorar casos y hallazgos similares.} 
El sistema permitirá seleccionar un caso o hallazgo para 
identificar otros elementos que compartan características o 
narrativas similares dentro del conjunto analizado. \textbf{[OD2]}

\textbf{T5. Consultar el documento original de un dato extraído.} 
El sistema permitirá acceder al fragmento del documento original 
para navegar por sus páginas y verificar el contexto exacto del 
que se extrajo cada dato. \textbf{[OD3]}

\section{Arquitectura del sistema}

\begin{figure}[htbp]
    \centering
    \includegraphics[width=\textwidth]{PIPElin.png} 
    \caption{Diseño de la arquitectura del sistema, organizado 
    en tres fases, extracción, agregación y exploración 
    interactiva.}
    \label{fig:Pipeline}
\end{figure}

Tal como se observa en la Figura~\ref{fig:Pipeline}, el sistema 
se organiza en tres fases. En la fase de \textbf{extracción}, 
cada informe se procesa de forma individual para identificar 
variables, actores, relaciones y ubicaciones, preservando la 
referencia al documento original. Los datos extraídos pasan a la 
fase de \textbf{agregación}, donde los informes de un mismo caso 
se agrupan, sus variables se sincronizan en el tiempo y se 
calculan las agregaciones geográficas. Finalmente, la fase de 
\textbf{exploración interactiva} expone esta representación al 
analista mediante vistas coordinadas.

\subsection{Extracción}

La fase de extracción procesa cada informe de forma individual 
y produce una representación estructurada de su contenido junto 
con anclajes precisos al documento original. Se compone de tres 
etapas, preprocesamiento, extracción semántica y mitigación de 
alucinaciones.

\subsubsection{Preprocesamiento}

El preprocesamiento convierte cada informe en una representación 
textual semiestructurada en la que cada elemento queda acompañado 
de su ubicación exacta en el documento, expresada como 
identificador, página y coordenadas. Optamos por una 
representación semiestructurada antes que por una estructura 
rígida porque los informes provienen de fuentes heterogéneas, y 
una plantilla fija habría requerido reglas específicas por tipo 
de documento, lo que comprometería la generalidad del 
sistema~\cite{cattan2021crossdoc}. Los informes pueden contener 
texto nativo o haber sido digitalizados mediante escaneo. En 
ambos casos utilizamos OpenDataLoader 
\ac{PDF}~\cite{opendataloader2024}, que aplica extracción directa al 
texto nativo y activa \ac{OCR} cuando lo detecta como necesario. A 
partir de esta salida se genera una representación simplificada 
que conserva solo el identificador y el texto de cada elemento, 
reduciendo el volumen de información enviada al modelo en la 
etapa siguiente.

\subsubsection{Extracción semántica}

La extracción semántica se realiza en dos consultas sucesivas 
al modelo. En la primera, se describe en lenguaje natural 
la información relevante del informe, incluyendo título, 
descripción breve, fechas, actores y ubicación geográfica. 
Para cada variable, clasificada según los seis tipos definidos 
en la Tabla~\ref{tab:tipos_variables}, el modelo devuelve tres 
elementos. Primero, el contenido de la variable. Segundo, un 
fragmento del texto que sirve de evidencia. Tercero, una 
justificación breve en lenguaje natural que explica cómo 
interpretó el fragmento para llegar a ese contenido. En la 
segunda consulta, el modelo organiza esa respuesta en la 
estructura de salida del sistema y conserva el identificador 
del elemento del que se tomó cada dato.

La decisión de exigir una justificación junto a cada dato 
responde a un argumento documentado en la 
literatura~\cite{yeh2025storyribbons}, donde se reporta que la 
capacidad de los \acs{LLM} para producir justificaciones de sus 
salidas ayuda a calibrar la confianza del usuario en la 
extracción. La justificación cumple un rol técnico adicional. 
Reduce el riesgo de respuestas no ancladas al texto, ya que 
obliga al modelo a articular el razonamiento como parte de la 
respuesta antes de comprometerse con un valor, en línea con el 
principio de razonamiento encadenado~\cite{wei2022}. Además, provee 
al analista un punto de partida para auditar el caso, sin 
tener que reconstruir desde cero por qué el modelo eligió un 
valor sobre otro.

La separación en dos consultas responde a una decisión 
adicional. Pedirle al modelo que razone sobre el contenido y 
al mismo tiempo lo serialice en un esquema rígido degrada la 
calidad de ambas tareas, según evidencia reportada en la 
literatura sobre extracción 
estructurada~\cite{li2024simple}. Al separar razonamiento y 
formato, la primera consulta se aproxima a un razonamiento 
encadenado, donde el modelo expresa libremente su análisis 
junto con la justificación. La segunda se limita a una tarea 
de transformación, donde el riesgo de alucinación es menor 
porque el contenido ya fue producido y solo resta darle 
estructura.

La tipología de seis tipos de variables que adoptamos refleja 
las distintas formas en que un dato puede evolucionar entre 
informes. Algunas variables cambian de estado, otras avanzan 
hacia una meta numérica, otras describen relaciones, y otras 
verifican una condición. Distinguirlas explícitamente permite 
que la fase de agregación trate cada tipo con la lógica 
apropiada y que la verificación elija la estrategia adecuada 
para cada dato.



\begin{table}[t]
\centering
\caption{Tipos de variables del modelo de datos.}
\label{tab:tipos_variables}
\renewcommand{\arraystretch}{1.5}
\begin{tabular}{@{}p{2.1cm} p{4.5cm} p{4cm} p{3.2cm}@{}}
\toprule
\textbf{Tipo} & \textbf{Descripción} & \textbf{Detalle interno} 
& \textbf{Ejemplo} \\
\midrule
\rowcolor{rowgray}
Seguimiento  & Datos que evolucionan entre informes, pudiendo 
resolverse, persistir o generar nuevas consecuencias & Texto, 
estado y fuente & Situaciones adversas detectadas \\
Progreso     & Valores numéricos que avanzan hacia una meta, 
comparados con un estimado por informe & Valor observado, meta, 
valor estimado y fuente & Porcentaje de presupuesto ejecutado \\
\rowcolor{rowgray}
Estado       & Categoría que describe la condición de algo en 
un informe específico & Valor categórico, nota y fuente & 
Tipo de acción (correctiva, preventiva) \\
Relacional   & Actores que participan en el caso y los roles 
que desempeñan & Lista de actores con nombre, rol y fuente & 
Organismos fiscalizados y supervisores \\
\rowcolor{rowgray}
Narrativo    & Texto estructurado con título y contenido 
extenso & Título, cuerpo del texto y fuente & Conclusiones y 
recomendaciones formales \\
Verificación & Condición booleana respaldada por evidencia y 
nivel de cumplimiento & Valor booleano, evidencia y fuente & 
Plazo de respuesta cumplido \\
\bottomrule
\end{tabular}
\end{table}

Cada dato extraído se almacena junto a tres anclajes que lo 
conectan con el documento original. La fuente, que indica de 
dónde proviene mediante el identificador del elemento, la 
página y las coordenadas. El fragmento de evidencia, que 
contiene el texto exacto desde el que se realizó la extracción. 
Y la justificación, que articula la interpretación del modelo 
sobre ese fragmento. Estos tres anclajes cumplen funciones 
complementarias. Permiten al analista consultar el documento 
desde la interfaz, habilitan la verificación automática de la 
etapa siguiente, y sirven de evidencia explícita para que el 
analista pueda calibrar su confianza en cada extracción.

\subsubsection{Mitigación de alucinaciones}

Dado que la alucinación en \acs{LLM} es una limitación innata que 
solo puede mitigarse y no eliminarse~\cite{xu2024hallucination}, 
el sistema adopta un esquema de verificación en capas que 
reduce su incidencia y preserva la trazabilidad necesaria para 
que el analista identifique los errores que escapen al sistema.

La verificación automática se aplica antes de pasar a la fase 
de agregación. La estrategia se adapta a la naturaleza del 
dato, distinguiendo entre valores literales y valores 
inferidos, ya que cada uno admite un tipo distinto de 
contraste contra la fuente.


Los valores literales, como fechas, nombres y montos, se 
verifican mediante coincidencia de texto contra el fragmento 
indicado por su fuente. Esta verificación detecta directamente 
las fabricaciones del modelo a un costo computacional mínimo, 
en línea con técnicas de auto-verificación 
factual~\cite{manakul2023selfcheckgpt}. Si el valor aparece 
en el fragmento, se considera verificado. En caso contrario, 
se marca como sospechoso y se envía a una nueva consulta al 
modelo, que recibe el dato junto al fragmento y debe 
reextraerlo o confirmar su ausencia. El resultado reemplaza 
al dato inicial.

Los valores inferidos, que el modelo deduce a partir del 
texto en lugar de tomarlos literalmente, no admiten una 
verificación por coincidencia. Para estos casos adoptamos un 
esquema de verificación encadenada 
(\textit{chain-of-verification})~\cite{dhuliawala2024chainofverification}, 
en el que el modelo revisa su propia respuesta contra el 
fragmento original y la justificación que produjo, y confirma 
si la inferencia se desprende razonablemente del texto. Si 
falla, el dato se reextrae añadiendo el contexto explícito.

La combinación de ambos mecanismos responde a un compromiso 
entre costo computacional y capacidad de razonamiento. La 
comparación de texto resuelve los casos simples sin invocar 
al modelo, lo que reduce el costo y elimina la posibilidad 
de alucinación en datos verificables sintácticamente. La 
verificación encadenada se reserva para los datos que 
requieren interpretación, donde el costo adicional se 
justifica por el riesgo de error. Esta selectividad reduce 
la tasa global de alucinaciones sin penalizar el tiempo de 
procesamiento.

Aun así, el resultado de la verificación automática no 
constituye una garantía. Por la limitación teórica antes 
señalada, una fracción de errores puede atravesar ambos 
filtros, especialmente en casos donde la evidencia textual es 
ambigua o donde el modelo refuerza su propia interpretación 
errónea durante la verificación encadenada. Por ello, los 
tres anclajes producidos durante la extracción semántica, 
fuente, fragmento de evidencia y justificación, se conservan 
en la representación final del informe. La interfaz expone 
estos anclajes en la fase de exploración, de modo que el 
analista pueda contrastar cada dato con su contexto original 
y emitir un juicio informado. Esta separación de 
responsabilidades, donde el sistema reduce la incidencia de 
errores y el analista valida el resultado, es coherente con 
los marcos de análisis visual que conciben la interpretación 
como una tarea conjunta entre el cómputo automático y el 
razonamiento humano~\cite{keim2008visual, sacha2014}.

\subsection{Agregación}

La fase de agregación toma los informes procesados 
individualmente en la etapa anterior y los consolida en una 
representación unificada por caso. El reto central de esta 
fase es la resolución de correferencia entre informes, ya que 
una misma variable puede aparecer descrita con texto diferente 
en distintos informes del mismo caso, y sin un mecanismo de 
resolución los datos no pueden integrarse en una secuencia 
coherente~\cite{cattan2021crossdoc}. La fase se compone de 
cuatro pasos que se ejecutan en orden, agrupación por caso, 
emparejamiento semántico de variables, verificación cruzada y 
construcción del vector de características.

\subsubsection{Agrupación por caso}

Los informes que comparten un mismo caso se agrupan, 
conformando su secuencia ordenada en el tiempo. Los datos 
generales del caso, como el título o la descripción, se fusionan 
en una representación única tomando como base el primer informe 
disponible y completando los campos ausentes con los informes 
posteriores. Esta política de fusión refleja el supuesto de que 
los informes posteriores complementan, antes que contradicen, 
la información del primero. Las contradicciones reales se 
detectan en el paso de verificación cruzada.

\subsubsection{Emparejamiento semántico de variables}

Una vez agrupados los informes de un caso, el sistema compara 
las variables entre informes consecutivos para identificar 
cuáles representan el mismo elemento descrito con texto 
diferente. Este emparejamiento se realiza mediante una consulta 
al \acs{LLM}, que recibe las variables de informes consecutivos junto 
con su justificación y determina cuáles corresponden al mismo 
elemento evolucionado, registrando el cambio observado entre 
informes.

La decisión de delegar el emparejamiento al modelo en lugar de 
basarlo únicamente en similitud léxica responde a la naturaleza 
del texto institucional. Una misma variable puede reformularse 
de un informe al siguiente, cambiar de tono o incorporar nuevos 
detalles, sin que ello implique un cambio en el referente. 
Métricas léxicas tradicionales, como similitud de cadenas o 
\textit{embeddings} de palabra, no capturan adecuadamente estos casos 
cuando la reformulación es sustantiva. Trabajos recientes en 
consistencia entre resúmenes han mostrado que delegar la 
comparación a un modelo de lenguaje produce decisiones más 
robustas frente a paráfrasis y reformulaciones~\cite{laban2022summac}, 
y es ese el tipo de robustez que el dominio requiere.

Durante este proceso, las variables que no aparecen en todos 
los informes se registran con valor nulo, distinguiendo 
explícitamente entre un dato ausente y un dato con valor cero. 
Esta distinción es necesaria porque, en un caso, la ausencia de 
una variable en un informe específico es una señal informativa 
en sí misma, no un error de extracción.

\subsubsection{Verificación cruzada}

Con el caso ya reconstruido, el sistema revisa inconsistencias 
entre informes que no eran visibles al procesar cada documento 
por separado. Esto incluye fechas contradictorias, nombres 
distintos para un mismo actor y hallazgos cuyo estado declarado 
no es consistente con su evolución a lo largo del 
caso~\cite{manakul2023selfcheckgpt}.

Este paso es complementario al de mitigación de alucinaciones 
de la fase anterior. Mientras aquella verificación contrasta 
cada dato contra su fragmento fuente, la verificación cruzada 
contrasta cada dato contra el resto de los informes del caso. 
Capturan tipos de error distintos. Una fecha extraída 
correctamente de un fragmento puede entrar en conflicto con la 
fecha extraída correctamente de otro fragmento si los documentos 
mismos no coinciden, y solo una verificación a nivel de caso 
detecta ese tipo de inconsistencia.
%¿por qué no usaron el LLM para resolver?", tu respuesta tiene que ser exactamente la del segundo párrafo. Lo intentamos así inicialmente, pero el modelo terminaba inventando razones para preferir uno u otro valor sin que existiera evidencia textual real para hacerlo. Eso es contradictorio con el objetivo de mitigar alucinaciones. La decisión de marcar y exponer es más conservadora, más honesta, y además aprovecha el valor informativo del conflicto.
Cuando se detecta una inconsistencia, el sistema no la resuelve 
automáticamente. Conserva los valores en conflicto junto con 
sus respectivos anclajes al documento fuente y los marca como 
inconsistentes en la representación final del caso. Esta 
decisión responde a dos consideraciones. Por un lado, una 
resolución automática mediante una nueva consulta al modelo 
introduce el mismo riesgo de alucinación que el sistema busca 
mitigar, ya que el modelo podría producir una respuesta 
plausible sin que existan elementos para preferir un valor 
sobre el otro. Por otro, una inconsistencia entre informes de 
un mismo caso es a menudo informativa en sí misma. Cuando dos 
documentos del mismo expediente reportan estados contradictorios 
sobre un mismo hallazgo, esa contradicción suele ser parte del 
fenómeno bajo análisis, no un error de extracción que conviene 
resolver. Marcar la inconsistencia y exponerla al analista 
preserva esa señal y respeta la división de roles entre el 
sistema y el analista propia del análisis 
visual~\cite{keim2008visual}.

Solo se aplica una normalización determinística para 
inconsistencias puramente sintácticas, como variantes de 
formato de fecha o diferencias de mayúsculas y acentos en 
nombres propios, donde la equivalencia es decidible sin 
interpretación.

\subsubsection{Construcción del vector de características}

Por último, el sistema calcula un vector de características 
para cada caso. Este vector combina dos componentes 
heterogéneos. El primero es un conjunto de métricas numéricas 
derivadas de las variables del caso, como la proporción de 
hallazgos de Seguimiento cerrados, el porcentaje de avance de 
las variables de Progreso, la proporción de condiciones de 
Verificación cumplidas y el número de actores involucrados. El 
segundo es un \textit{embedding} semántico generado a partir 
de las conclusiones registradas en los informes del caso. Para 
el \textit{embedding} se utiliza un modelo de \textit{Sentence 
Transformers}~\cite{reimers2019sbert}, que proyecta el texto a 
un espacio vectorial denso preservando su significado semántico.

La decisión de combinar componentes numéricos y semánticos en 
un solo vector responde al objetivo de soportar las dos 
formas en que un analista puede juzgar la similitud entre 
casos. Dos casos pueden parecerse porque sus indicadores 
cuantitativos coinciden, o porque sus conclusiones narrativas 
describen situaciones análogas. Un vector que ignore una de 
las dos dimensiones colapsaría una distinción que el analista 
sí percibe. La combinación obliga, no obstante, a normalizar 
ambos componentes antes de concatenarlos para que ninguno 
domine la métrica de distancia, ajuste que se realiza mediante 
escalado por desviación estándar.

El vector resultante es el insumo de la vista de similitud, 
que lo proyecta a dos dimensiones mediante 
\ac{UMAP}~\cite{mcinnes2018umap}. Optamos por \acs{UMAP} frente a 
alternativas como t-SNE porque preserva mejor la estructura 
global del conjunto de datos, lo que es relevante para una 
tarea de exploración donde la posición relativa entre grupos 
de casos comunica información al analista, no solo la 
vecindad local entre puntos individuales.

Al finalizar esta fase, el sistema produce tres archivos por 
conjunto de casos analizado. Un consolidado por caso con todos 
sus informes y variables sincronizadas. Un archivo de 
coordenadas geográficas por caso. Y un archivo de vectores de 
características, uno por caso. Estos archivos constituyen la 
entrada de la fase de exploración interactiva.

\subsection{Exploración interactiva}

La fase de exploración interactiva expone la representación 
unificada producida por las fases anteriores al analista. Su 
diseño parte del principio de vistas múltiples coordinadas, 
que sostiene que un fenómeno multidimensional se comprende 
mejor cuando cada dimensión se representa con la 
visualización más adecuada y las vistas se vinculan mediante 
interacción 
sincronizada~\cite{wang2000guidelines, roberts2007cmv}.

\begin{figure}[htbp]
    \centering
    \includegraphics[width=\textwidth]{vista_p.png}
    \caption{Interfaz general del sistema sobre un caso de 
    estudio del Proceso Electoral 2026, donde cuatro Oficinas 
    Descentralizadas de Procesos Electorales se analizan en 
    paralelo. La interfaz se compone de cuatro regiones 
    coordinadas, la vista geoespacial (\textbf{A}), la vista 
    temporal (\textbf{B}), la vista de similitud (\textbf{C}) 
    y la vista de detalle (\textbf{D}).}
    \label{fig:vista_principal}
\end{figure}

La exploración sigue un flujo de cuatro etapas. El analista 
parte de una vista panorámica del conjunto de casos, ya sea 
sobre el territorio en la vista geoespacial (\textbf{A}) o 
sobre el espacio semántico de similitud (\textbf{B}). Al 
seleccionar uno o varios casos en cualquiera de las dos, la 
selección se propaga a la otra y se construye la vista 
temporal (\textbf{C}) sobre los casos elegidos. Desde la 
vista temporal, el analista puede inspeccionar un informe, 
una variable o una región de la línea de tiempo, lo que 
abre la vista de detalle (\textbf{D}) con la información 
asociada. Y desde la vista de detalle, los enlaces a los 
fragmentos fuente abren la vista de evidencia, que muestra 
el documento original con el fragmento resaltado. Un panel 
de filtros adaptativos en el encabezado complementa este 
flujo, permitiendo restringir la exploración por tiempo, 
espacio y características del caso.

A continuación, detallamos cada componente siguiendo el recorrido del usuario: desde las primeras vistas de exploración hasta el acceso a la evidencia original.

\subsubsection{Vista geoespacial}

La vista geoespacial (\textbf{A} en la 
Figura~\ref{fig:vista_principal})~localiza los casos sobre 
un mapa interactivo y soporta la tarea T3. Aprovecha que 
cada caso tiene una ubicación asociada, calculada en la 
fase de agregación a partir de las referencias geográficas 
extraídas de los informes.

La vista geoespacial soporta selección de 
casos individuales o de regiones territoriales. La 
selección se propaga a la vista de similitud, donde los 
puntos correspondientes se resaltan, y a la vista 
temporal, que se construye sobre el subconjunto elegido.


\subsubsection{Vista de similitud}

La vista de similitud (\textbf{B} en la 
Figura~\ref{fig:vista_principal})~proyecta el conjunto 
completo de casos en un espacio bidimensional para 
soportar la tarea T4, donde el analista busca identificar 
grupos de casos parecidos o casos atípicos respecto al 
conjunto. Su fuente es el vector de características 
construido en la fase de agregación, que combina 
indicadores numéricos con un \textit{embedding} semántico 
de las conclusiones del caso.

La proyección se realiza con \acs{UMAP}~\cite{mcinnes2018umap}, 
cuya elección frente a alternativas como t-SNE ya fue 
discutida en la fase de agregación. La representación es 
un gráfico de dispersión donde cada punto corresponde a 
un caso. La cercanía espacial entre puntos indica que los 
casos comparten un perfil cuantitativo y narrativo 
similar, lo cual permite al analista detectar agrupaciones 
por su sola geometría sin tener que inspeccionarlos uno 
por uno.

La vista de similitud es una entrada complementaria a la 
geoespacial. Mientras la geoespacial agrupa casos por 
proximidad territorial, la de similitud los agrupa por 
proximidad semántica y cuantitativa. Esta dualidad permite 
que el analista descubra, por ejemplo, que dos casos 
geográficamente distantes presentan trayectorias casi 
idénticas, o que dos casos vecinos difieren en su 
comportamiento. Al posicionar el cursor sobre un punto, 
una ventana flotante muestra el resumen del caso. La 
selección de un punto, o de un grupo de puntos mediante 
\textit{lasso}, propaga el contexto a la vista 
geoespacial y construye la vista temporal sobre los casos 
elegidos.

\subsubsection{Filtros adaptativos}

El encabezado de la interfaz expone un panel de filtros 
que opera sobre la selección activa. Tres dimensiones de 
filtrado están disponibles. La dimensión temporal, 
mediante un control deslizante que limita la ventana de 
fechas considerada. La dimensión espacial, mediante 
selectores territoriales que restringen la región 
analizada. Y la dimensión de características, mediante 
controles que filtran por atributos de las variables, 
como el estado de una variable de Seguimiento o el rango 
de valores de una variable de Progreso.

Los controles de la dimensión de características se 
adaptan al tipo de variable activa. Esta adaptación responde a un 
principio elemental de diseño de interfaz, según el cual 
los controles deben corresponder a la estructura de los 
datos sobre los que 
operan~\cite{shneiderman1996tasks}, y evita el patrón 
común de exponer controles que no aplican a la variable 
activa, que confunde al usuario y degrada la experiencia 
de exploración.

\subsubsection{Vista temporal}

La vista temporal (\textbf{C} en la 
Figura~\ref{fig:vista_principal})~se construye sobre los 
casos seleccionados en las vistas de entrada y soporta 
las tareas T1 y T2. Su representación se adapta al tipo 
de variable bajo análisis, siguiendo el principio de que 
una sola codificación visual no puede expresar 
adecuadamente la heterogeneidad de los seis tipos del 
modelo de datos.

Por defecto, cuando ninguna variable específica está 
seleccionada, la vista temporal muestra los informes de 
los casos como un diagrama de Gantt, donde el eje 
horizontal corresponde al tiempo y el eje vertical 
organiza los casos en filas paralelas. Esta 
representación responde a una propiedad estructural del 
dominio. Los informes institucionales se caracterizan 
por intervalos con duración explícita, como un plazo de 
respuesta o una ventana de ejecución, intercalados con 
eventos puntuales, como la emisión del informe mismo. 
La sintaxis del Gantt preserva ambas clases y permite 
identificar visualmente retrasos, periodos sin actividad 
y solapamientos entre casos. La organización en filas 
paralelas, además, soporta la comparación directa entre 
casos sin requerir una segunda vista.

Al seleccionar una variable de \textbf{Seguimiento}, la 
vista temporal cambia a un diagrama de Sankey, donde 
cada flujo representa la transición de un hallazgo entre 
estados a lo largo de los informes. Esta elección 
responde a la estructura semántica del tipo. Las 
variables de Seguimiento describen elementos que se 
resuelven, persisten o generan nuevas consecuencias 
entre informes, lo cual es exactamente el patrón que 
los diagramas de Sankey representan, flujos entre 
estados con grosor proporcional al volumen del 
flujo~\cite{riehmann2005sankey}. La representación 
permite que el analista identifique, por ejemplo, qué 
proporción de hallazgos abiertos en un informe inicial 
quedó cerrada en un informe posterior, y qué proporción 
persistió o derivó en nuevos hallazgos.

Al seleccionar una variable de \textbf{Progreso}, la 
vista temporal cambia a un diagrama de líneas, donde una 
línea representa el valor observado a lo largo del 
tiempo y una segunda línea, en estilo discontinuo, 
representa el valor estimado o meta. Esta es la 
codificación canónica para series temporales numéricas 
con referencia~\cite{cleveland1985elements} y permite 
detectar visualmente desvíos respecto a la meta sin 
recurrir a cálculos auxiliares.

Para los cuatro tipos restantes, la vista temporal adopta 
codificaciones diferenciadas según la estructura semántica 
de cada uno. Una variable de \textbf{Estado}, que asigna 
una categoría a cada informe, se representa como una 
banda horizontal segmentada por color, siguiendo el 
patrón de las visualizaciones de tipo \textit{swimlane} 
que codifican secuencias categóricas a lo 
largo del tiempo~\cite{wongsuphasawat2011lifeflow}. Una 
variable \textbf{Relacional}, que describe los actores 
involucrados y sus roles, se representa como un grafo 
dinámico cuando el conjunto de actores es pequeño, o 
como una matriz de adyacencia ordenada cronológicamente 
cuando es 
denso~\cite{beck2017taxonomy}. Una variable de tipo 
\textbf{Verificación}, que registra el cumplimiento 
booleano de una condición a lo largo de los informes, se 
representa como un \textit{dot plot} binario donde cada 
fila corresponde a un caso y cada columna a un informe, 
codificación que permite identificar rápidamente patrones 
de incumplimiento sistemático. Las variables de tipo 
\textbf{Narrativo}, finalmente, no admiten una 
codificación temporal compacta porque su contenido es 
texto extenso. En su lugar, se exponen como tarjetas en 
la vista de detalle al seleccionar el informe 
correspondiente, manteniendo el acceso al texto sin 
forzar una representación visual que perdería 
información.

\subsubsection{Vista de detalle}

La vista de detalle (\textbf{D} en la 
Figura~\ref{fig:vista_principal})~se activa al 
seleccionar un informe, un caso o una región temporal en 
las otras vistas. Presenta la información asociada a esa 
selección, organizada en bloques que cubren los 
metadatos del elemento, las variables registradas en él 
y los enlaces a sus fragmentos fuente.

Para un informe seleccionado, la vista muestra el título 
del hito, las fechas relevantes (emisión del informe, 
plazo de respuesta), las situaciones identificadas y la 
referencia al documento original. Cada dato visible en 
esta vista es interactivo y abre la vista de evidencia 
sobre el fragmento correspondiente del documento.

Esta vista funciona como puente entre las 
representaciones agregadas y la evidencia textual. Las 
vistas A, B y C operan sobre representaciones derivadas 
del proceso de extracción, y aunque son útiles para 
reconocer patrones, no exponen el texto original que 
sustenta cada dato. La vista de detalle traduce un 
elemento seleccionado en sus componentes verificables, y 
expone los enlaces que permiten descender al documento 
fuente.

\subsubsection{Vista de evidencia}

La vista de evidencia materializa los anclajes producidos 
durante la extracción y soporta la tarea T5. Cuando el 
analista activa cualquier enlace desde la vista de 
detalle, la vista de evidencia se abre como una ventana 
embebida sobre la interfaz, carga el documento original, 
navega a la página correspondiente y resalta el 
fragmento exacto del que se obtuvo el dato. La 
Figura~\ref{fig:vista_pdf} muestra esta interacción 
sobre un informe de auditoría, donde el recuadro 
amarillo señala el fragmento fuente del dato 
seleccionado.

\begin{figure}[htbp]
    \centering
    \includegraphics[width=\textwidth]{vista_pdf.png}
    \caption{Vista de evidencia. El sistema abre el 
    documento original, navega a la página correspondiente 
    y resalta el fragmento desde el cual se extrajo el 
    dato seleccionado.}
    \label{fig:vista_pdf}
\end{figure}

La vista ofrece tres elementos de soporte para la 
verificación. Primero, la navegación entre páginas del 
documento permite consultar el contexto previo y 
posterior al fragmento, ya que un fragmento aislado 
puede no ser suficiente para juzgar la corrección de la 
extracción. Segundo, el resaltado visual del fragmento 
fuente reduce el costo cognitivo de localizarlo dentro 
de la página, especialmente en informes densos donde el 
fragmento relevante puede ser una sola línea entre 
cientos. Tercero, junto al fragmento se expone la 
justificación que el modelo articuló al realizar la 
extracción, lo que permite al analista contrastar no 
solo el dato sino también el razonamiento que lo 
produjo. Esta vista es central al diseño del sistema. Las fases 
anteriores adoptaron una posición explícita sobre la 
limitación teórica de los \acs{LLM} frente a la 
alucinación, asumiendo que el sistema mitiga pero no 
elimina los errores. Esa posición sería declarativa sin 
un mecanismo concreto que permita al analista verificar 
cada dato. La vista de evidencia es ese mecanismo. Al 
colocar el documento original a un clic de cualquier 
dato extraído, el sistema se alinea con la noción de calibración 
de confianza reportada en la literatura~\cite{yeh2025storyribbons}.
 %Inserta el capítulo 4
\input{conclusiones} %Inserta el capítulo 5

%%%%%%%%%%%%%%%%%%%%%%%%%%%%%%%%%%%%%%%%%%%%%%%%%%%%%%%%%%%%%%%%%%%%%%

\bibliographystyle{apalike}

\bibliography{Bibliog}
\addcontentsline{toc}{chapter}{Bibliografía}
\end{document}
